%% 9장 측도론적 확률 (새 목차, 2026-11-25)
\chapter{Measure-Theoretic Probability}
\label{ch:probability}

\abstract{확률론은 전체 측도가 $1$인 측도론에 독립과 조건부 기댓값을 더한 것이다. 이 장은 분포와 분위수 함수, 기댓값, 독립, 조건부 기댓값, 수렴의 양식, 특성함수, 큰 수의 법칙과 중심극한정리, 경험 분포 함수를 세운다. 구현으로는 stardust의 이산 분위수 탐색과 그 후방 오차 계약, 역변환 표집, 꼬리 확률과 로그 확률의 계산을 따라가고 극한 정리의 속도를 정확한 계산으로 잰다.}

통계 계산의 거의 모든 출력은 확률이다. p값, 신뢰구간의 끝점, 표본의 분위수, 난수, 몬테카를로 추정치가 모두 그렇다. 이런 값이 무엇을 뜻하고 얼마나 정확해야 하는지를 말하려면 확률을 측도로 세워야 한다. 측도론 위에서는 분포 함수와 분위수 함수가 정확히 무엇인지, 표본평균이 어떤 뜻에서 참평균으로 가는지, 그 속도가 얼마인지를 다룰 수 있다. 이 장에서는 그 바탕을 다진 뒤 stardust가 분포 함수와 분위수 함수, 난수를 계산하는 방식을 이론에 맞대어 본다. 독자가 측도론을 안다고 보고 그 정리(단조·지배 수렴, Fubini, Radon--Nikodym)는 증명 없이 쓴다.

\section{Probability Spaces}
\label{sec:probability-space}

\term{확률 공간}{probability space}은 측도 공간 $(\Omega,\mathcal F,P)$로서 $P(\Omega)=1$인 것이다. $\Omega$의 원소를 \term{표본}{outcome}, $\mathcal F$의 원소를 \term{사건}{event}이라 부른다. 측도의 성질이 그대로 확률의 성질이 된다. 서로소인 사건들의 확률은 더해진다. 늘어나는 사건열 $A_n\uparrow A$에서 $P(A_n)\uparrow P(A)$이고 유한 측도이므로 줄어드는 사건열 $A_n\downarrow A$에서도 $P(A_n)\downarrow P(A)$이다. 이 공리 체계는 Kolmogorov가 1933년에 세웠다\cite{kolmogorov1933}.

측도가 같은지 확인할 때는 생성족에서만 보면 된다. 교집합에 닫힌 집합족을 $\pi$-체계, 전체를 담고 $A\subseteq B$인 두 원소의 차집합 $B\setminus A$와 늘어나는 합집합에 닫힌 집합족을 $\lambda$-체계라 한다. Dynkin의 $\pi$--$\lambda$ 정리에 따르면 $\pi$-체계를 담는 $\lambda$-체계는 그 $\pi$-체계가 생성한 $\sigma$-대수를 담는다. 증명의 줄거리는 이렇다. $\pi$-체계 $\mathcal P$를 담는 가장 작은 $\lambda$-체계 $\lambda(\mathcal P)$가 교집합에 닫혀 있음을 두 단계로 보인다. 먼저 $A\in\mathcal P$를 고정하면 $\{B:A\cap B\in\lambda(\mathcal P)\}$가 $\mathcal P$를 담는 $\lambda$-체계이므로 $\lambda(\mathcal P)$를 담고 이어 $B\in\lambda(\mathcal P)$를 고정해 같은 논법을 다시 쓴다. 교집합에 닫힌 $\lambda$-체계는 $\sigma$-대수이므로 끝난다(Billingsley\cite{billingsley1995} 정리 3.2). 그러면 두 확률 측도가 $\pi$-체계 위에서 같을 때 그것이 생성한 $\sigma$-대수 위에서도 같다($\{A:P(A)=Q(A)\}$가 $\lambda$-체계이기 때문이다). 이 장에서는 분포의 유일성을 보일 때도, 독립을 판정할 때도 이 원리를 쓴다.

사건열이 ``무한히 자주'' 일어나는지 판정하는 보조정리는 확률론에만 있는 도구이다. $\{A_n\text{ i.o.}\}=\bigcap_{m\ge1}\bigcup_{n\ge m}A_n$으로 쓴다.

\begin{lemma}[Borel--Cantelli]
\label{lem:borel-cantelli}
\begin{enumerate}
\item $\sum_nP(A_n)<\infty$이면 $P(A_n\text{ i.o.})=0$이다.
\item $A_n$들이 독립이고 $\sum_nP(A_n)=\infty$이면 $P(A_n\text{ i.o.})=1$이다.
\end{enumerate}
\end{lemma}

\begin{proof}
(i) 모든 $m$에서 $P(A_n\text{ i.o.})\le P(\bigcup_{n\ge m}A_n)\le\sum_{n\ge m}P(A_n)$이고 오른쪽은 수렴하는 급수의 꼬리이므로 $m\to\infty$에서 $0$이다. (ii) 여사건을 보면 된다. 독립이므로 $P(\bigcap_{n=m}^MA_n^c)=\prod_{n=m}^M(1-P(A_n))\le\exp(-\sum_{n=m}^MP(A_n))$이고($1-x\le e^{-x}$) $M\to\infty$에서 $0$이다. 따라서 모든 $m$에서 $P(\bigcap_{n\ge m}A_n^c)=0$이고 가산 합집합 $\bigcup_m\bigcap_{n\ge m}A_n^c=\{A_n\text{ i.o.}\}^c$의 확률도 $0$이다.
\end{proof}

두 반은 같은 경계의 양쪽을 다룬다. 독립인 균등 난수 $U_n$으로 $A_n=\{U_n\le1/n\}$과 $B_n=\{U_n\le1/n^2\}$을 만들면 $\sum1/n=\infty$이므로 $A_n$은 무한히 자주 일어나고 $\sum1/n^2<\infty$이므로 $B_n$은 유한 번만 일어난다. \cref{fig:borel-cantelli}에서 $A_n$의 누적 횟수는 기댓값 $\sum_{k\le n}1/k\approx\log n+0.577$을 따라 끝없이 늘고 $B_n$은 처음 몇 번 뒤로 멈춘다.

\begin{figure}[htbp]
\centering
\begin{tikzpicture}
\begin{semilogxaxis}[sdplot,xlabel={$n$},ylabel={$n$까지 일어난 횟수},xmin=1,xmax=1e6,ymin=0,ymax=20,width=0.7\textwidth,height=0.4\textwidth]
\addplot[black,const plot,thick] table[x=n,y=inv]{figures/09/bc.dat};
\addlegendentry{$A_k=\{U_k\le1/k\}$}
\addplot[black,const plot,dashed,thick] table[x=n,y=inv2]{figures/09/bc.dat};
\addlegendentry{$B_k=\{U_k\le1/k^2\}$}
\addplot[black!50,dotted,thick] table[x=n,y=logn]{figures/09/bc.dat};
\addlegendentry{$\log n+0.577$}
\end{semilogxaxis}
\end{tikzpicture}
\caption{Borel--Cantelli의 두 경우(같은 균등 난수 $10^6$개, 따로 빌드한 stardust의 \code{random}; \code{figures/09/extras.py}). $10^6$번 뒤 $A$는 $17$번(기댓값 $14.4$), $B$는 $1$번 일어났다.}
\label{fig:borel-cantelli}
\end{figure}

\section{Distributions and Quantile Functions}
\label{sec:distributions}

\subsection{Distribution Functions}

가측 함수 $X:(\Omega,\mathcal F)\to(\mathbb R,\mathcal B)$를 \term{확률변수}{random variable}라 한다($\mathcal B$는 Borel 집합족). $X$의 \term{분포}{distribution}는 밀어낸 측도 $P_X(B)=P(X\in B)$이고 $\mathbb R$ 위의 확률 측도이다. 통계에서는 거의 언제나 분포만 다루고 $\Omega$는 여러 확률변수를 한 곳에 함께 놓으려고 두는 공간일 뿐이다. 분포는
\[
F(x)=P(X\le x)=P_X\bigl((-\infty,x]\bigr)
\]
인 \term{누적분포함수}{cumulative distribution function}로 요약된다. $F$는 단조 증가하고 오른쪽 연속이며($(-\infty,x+1/n]\downarrow(-\infty,x]$에 연속성을 쓴다) $F(-\infty)=0$, $F(\infty)=1$이다. 반직선 $(-\infty,x]$들은 $\pi$-체계이고 $\mathcal B$를 생성하므로 분포는 누적분포함수로 유일하게 정해진다. 거꾸로 세 성질을 가진 함수는 모두 어떤 분포의 누적분포함수이고 이것은 다음 소절에서 구성으로 보인다. 원자 $x$에서는 $P(X=x)=F(x)-F(x-)$이다.

\subsection{The Quantile Function}

누적분포함수는 엄격히 증가하지 않을 수도 있고(밀도가 $0$인 구간) 연속이 아닐 수도 있으므로(원자) 역함수를 다음처럼 일반화한다.

\begin{definition}
\label{def:quantile}
누적분포함수 $F$의 \term{분위수 함수}{quantile function}를 $p\in(0,1)$에서
\[
F^{-}(p)=\inf\{x\in\mathbb R:F(x)\ge p\}
\]
로 정의한다.
\end{definition}

$x\to\infty$에서 $F(x)\to1>p$이고 $x\to-\infty$에서 $F(x)\to0<p$이므로 집합은 공집합이 아니고 아래로 유계여서 $F^{-}(p)$는 실수이다(\cref{fig:quantile}). 이 절의 논증은 모두 다음 동치를 쓴다.

\begin{lemma}
\label{lem:quantile-galois}
$p\in(0,1)$과 $x\in\mathbb R$에 대해 $F^{-}(p)\le x$인 것은 $p\le F(x)$인 것과 동치이다.
\end{lemma}

\begin{proof}
$p\le F(x)$이면 $x$가 정의의 집합에 들므로 $F^{-}(p)\le x$이다. 거꾸로 $F^{-}(p)\le x$라 하자. 하한의 정의에서 $F(x_n)\ge p$이고 $x_n\downarrow F^{-}(p)$인 열이 있다. $F$가 오른쪽 연속이므로 $F(F^{-}(p))=\lim F(x_n)\ge p$이고 $F$가 단조이므로 $F(x)\ge F(F^{-}(p))\ge p$이다.
\end{proof}

\begin{theorem}[역변환과 확률 적분 변환]
\label{thm:inverse-transform}
$U$가 $(0,1)$의 균등 분포를 따르면 $X=F^{-}(U)$의 누적분포함수는 $F$이다. 또 $F$가 연속이면 $F(X)$는 $(0,1)$의 균등 분포를 따른다.
\end{theorem}

\begin{proof}
\cref{lem:quantile-galois}에서 $P(X\le x)=P(F^{-}(U)\le x)=P(U\le F(x))=F(x)$이다. 둘째 주장: 같은 보조정리에 $x=X$를 넣으면 사건 $\{F(X)\ge p\}$와 $\{X\ge F^{-}(p)\}$가 같다. $F$가 연속이면 $(0,1)$의 모든 $p$가 $F$의 값이므로 $F(F^{-}(p))=p$이고(보조정리의 증명에서 $\ge p$이고 연속성과 하한의 정의에서 $\le p$이다) $X$에 원자가 없으므로 $P(X\ge y)=1-F(y)$이다. 따라서 $P(F(X)\ge p)=1-F(F^{-}(p))=1-p$이다.
\end{proof}

첫째 주장은 누적분포함수의 세 성질을 가진 $F$에서 분포를 구성하는 방법이기도 하다. $(0,1)$ 위의 Lebesgue 측도를 $F^{-}$로 밀어내면 된다. \cref{sec:inversion}의 \term{역변환 표집}{inverse transform sampling}은 이 주장을 그대로 구현한다. 둘째 주장을 \term{확률 적분 변환}{probability integral transform}이라 한다. 검정 통계량의 분포가 연속이면 귀무가설 아래 p값이 균등 분포를 따르는 것도 이 변환에서 나온다(11권). Lebesgue 분해 정리에서 확률 측도는 원자에 실린 이산 부분, Lebesgue 측도에 절대 연속인 부분(밀도 $f=dP_{ac}/d\lambda$), 원자가 없지만 Lebesgue 측도와 서로 특이한 부분의 합으로 유일하게 나뉜다. 마지막 부분의 예로 Cantor 함수를 누적분포함수로 갖는 분포가 있지만 통계 모형에서 실제로 만나는 분포는 밀도를 가진 분포와 정수에 실린 분포이고 가끔 그 둘의 혼합이 나온다.

확률변수를 변환하면 밀도도 바뀐다. $X$가 열린 집합 $U\subseteq\mathbb R^d$ 위에서 밀도 $f_X$를 갖고 $g:U\to V$가 연속 미분 가능한 전단사이며 야코비 행렬식이 $0$이 아니면 $Y=g(X)$의 밀도는
\[
f_Y(y)=f_X(g^{-1}(y))\,\bigl|\det Dg^{-1}(y)\bigr|
\]
이다. 다변수 치환 적분 공식에서
\[
P(Y\in B)=\int_{g^{-1}(B)}f_X\,dx=\int_Bf_X(g^{-1}(y))\,|\det Dg^{-1}(y)|\,dy
\]
이기 때문이다. 표집 공식의 상당수가 이 변환에서 나온다. 지수 표본 $E$에서 $-\log E$는 Gumbel, $E^{1/k}$는 Weibull, $e^{E/a}$는 Pareto가 되고(\cref{tab:samplers}) 독립인 정규 둘을 극좌표로 바꾸면 Box--Muller 방법이 된다.

\begin{figure}[htbp]
\centering
\begin{tikzpicture}
\begin{axis}[sdplot,width=0.45\textwidth,height=0.4\textwidth,xlabel={$x$},ylabel={$F(x)$},xmin=-0.5,xmax=4.5,ymin=0,ymax=1.05,title={누적분포함수}]
\addplot[black,thick,domain=-0.5:1,samples=2] {0};
\addplot[black,thick,domain=1:2,samples=40] {0.3*(x-1)};
\addplot[black,thick,domain=2:3,samples=2] {0.6};
\addplot[black,thick,domain=3:4.5,samples=40] {1-0.25*exp(-2*(x-3))};
\addplot[black,dashed] coordinates {(2,0.3) (2,0.6)};
\addplot[black,dashed] coordinates {(3,0.6) (3,0.75)};
\addplot[only marks,mark=*,mark size=1.3pt] coordinates {(2,0.6) (3,0.75)};
\addplot[only marks,mark=o,mark size=1.3pt] coordinates {(2,0.3) (3,0.6)};
\end{axis}
\end{tikzpicture}\hfill
\begin{tikzpicture}
\begin{axis}[sdplot,width=0.45\textwidth,height=0.4\textwidth,xlabel={$p$},ylabel={$F^{-}(p)$},xmin=0,xmax=1,ymin=-0.5,ymax=4.5,title={분위수 함수}]
\addplot[black,thick,domain=0.001:0.3,samples=40] {1+x/0.3};
\addplot[black,thick] coordinates {(0.3,2) (0.6,2)};
\addplot[black,thick] coordinates {(0.6,3) (0.75,3)};
\addplot[black,thick,domain=0.75:0.995,samples=60] {3-0.5*ln((1-x)/0.25)};
\addplot[only marks,mark=*,mark size=1.3pt] coordinates {(0.6,2) (0.75,3)};
\addplot[black,dashed] coordinates {(0.6,2) (0.6,3)};
\end{axis}
\end{tikzpicture}
\caption{원자가 둘($x=2$에 $0.3$, $x=3$에 $0.15$)이고 $[2,3)$에서 평평한 누적분포함수(왼쪽)와 그 분위수 함수(오른쪽). 원자는 분위수 함수의 평평한 구간이 되고 평평한 구간은 분위수 함수의 뜀이 된다. $F^{-}$는 왼쪽 연속이다.}
\label{fig:quantile}
\end{figure}

\section{Computing Quantiles}
\label{sec:computing-quantiles}

\subsection{What a Computed Quantile Can Promise}

연속 분포의 분위수는 $F(x)=p$의 근이므로 4장의 조건수 언어로 정확도를 말할 수 있다. $p$를 상대 $\delta$만큼 움직이면 근은 $\Delta x\approx p\delta/f(x)$만큼 움직이므로 상대 조건수는
\[
\kappa(p)=\left|\frac{p}{x\,f(x)}\right|,\qquad x=F^{-}(p)
\]
이다. 꼬리로 가면 $f(x)$가 작아져도 $p$도 함께 작아지므로 정규 분포의 아래 꼬리에서 $\kappa\approx1/x^2$로 오히려 잘 조건화되어 있다. 반대로 $p$를 큰 쪽에서 받으면 $p\approx1$이 상대 $u$로 반올림된 값이라 $1-p$의 정보가 이미 사라졌을 수 있다. 이 때문에 위쪽 꼬리의 분위수는 작은 수 $q=1-p$를 받는 역생존 함수 $\bar F^{-}(q)$로 따로 둔다. stardust의 분포 객체는 \code{icdf}, \code{iccdf}와 그 로그판 \code{ilogcdf}, \code{ilogccdf}를 두고 연속 분포에서는 $|\Delta x/x|\le10^{-14}\max(1,\kappa)$를 약속한다.

이산 분포에서는 사정이 근본적으로 다르다. 정수에 실린 분포에서 $F^{-}(p)$는 $F(k)\ge p$인 가장 작은 정수 $k$이므로 $p$가 어떤 원자의 누적 확률 $F(k)$와 정확히 같으면 그 원자가 답이고 $p$가 조금만 커져도 다음 원자가 답이 된다. 분위수 함수가 $p=F(k)$마다 불연속이므로 조건수는 무한대이다. 평균 $3$인 Poisson 분포에서 $F(2)=e^{-3}\cdot8.5=0.42319008112684351\ldots$이고 그에 가장 가까운 binary64 수 $0.42319008112684353$은 참값보다 $1.6\times10^{-17}$ 크다. 그 수의 정확한 분위수는 $3$이지만 바로 아래 이웃 double의 분위수는 $2$이다. 계산한 누적 확률에 반올림이 섞여 있는 한 원자 근처의 $p$에서 답이 한 칸 어긋나는 일은 막을 수 없다.

\begin{definition}[후방 오차 계약]
\label{def:backward-quantile}
이산 분포의 계산한 분위수 $\hat k$가 \term{후방 오차}{backward error} $\delta$ 안에서 옳다는 것은 $|\tilde p-p|\le\delta\min(p,1-p)$인 어떤 $\tilde p$에 대해 $\hat k=F^{-}(\tilde p)$라는 것이다.
\end{definition}

stardust는 $\delta=10^{-14}$의 후방 오차를 약속한다(설계 문서 T25). 작은 쪽 $\min(p,1-p)$에 상대적이므로 위쪽 꼬리에서도 계약의 뜻이 살아 있다. 위의 Poisson 예에서 따로 빌드한 stardust의 \code{Poisson(mu=3).icdf}는 $\mathrm{fl}(F(2))$에서 $2$를 돌려준다. 정확한 분위수 $3$과는 다르지만 $p$를 $10^{-16}$보다 작게 움직이면 $2$가 정확한 답이므로 계약 안이다(\code{figures/09/tails\_sd.py}). 이렇게 답하는 데에는 이유도 있다. 사용자는 대개 $p$를 \code{cdf(2)}로 계산해 넣으므로 그때 $2$를 돌려주어야 \code{icdf(cdf(k))==k}가 성립한다.

\subsection{Searching for a Discrete Quantile}

정수 $k$에서 누적 확률을 계산할 수 있으면 분위수는 단조인 술어 $F(k)\ge p$가 처음 참이 되는 정수이다. 탐색 범위가 무한할 수 있으므로(Poisson, 음이항) stardust는 근사로 얻은 시작점에서 보폭을 $1,2,4,\dots$로 두 배씩 늘려 술어가 바뀌는 구간을 찾고(\term{갤로핑}{galloping}) 그 구간을 이분한다(\cref{alg:discrete-icdf}). 시작점에서 답까지의 거리를 $\ell$이라 하면 술어 계산은 $O(\log\ell)$번이다.

\begin{algorithm}[htbp]
\begin{algorithmic}[1]
\Require 술어 $\mathrm{pred}$(단조, 정수에서), 지지 $[k_{\mathrm{lo}},k_{\mathrm{hi}}]$, 시작 추측 $g$
\Ensure $\mathrm{pred}(k)$가 참인 가장 작은 $k$
\State $g\gets\lfloor g\rfloor$를 지지로 자른다
\If{$\mathrm{pred}(g)$} \Comment{답은 $g$ 이하: 아래로 갤로핑}
  \State $b\gets g$; $a\gets k_{\mathrm{lo}}-1$; 보폭 $1,2,4,\dots$로 $c=\max(k_{\mathrm{lo}},b-\text{보폭})$을 보며($b>k_{\mathrm{lo}}$인 동안) $\mathrm{pred}(c)$가 거짓이면 $a\gets c$, 멈춤, 참이면 $b\gets c$
\Else \Comment{답은 $g$ 위: 위로 갤로핑}
  \State $a\gets g$; $b\gets k_{\mathrm{hi}}$; 보폭을 두 배씩 늘려 $c=a+\text{보폭}$을 본다($c\ge k_{\mathrm{hi}}$이면 멈춘다)
  \State \textbf{if} $c\ge2^{1023}$ \textbf{then} 가장 큰 double에서 $\mathrm{pred}$를 보고 거짓이면 \Return $\infty$, 참이면 $b\gets$ 그 값
  \State $\mathrm{pred}(c)$가 참이면 $b\gets c$, 아니면 $a\gets c$
\EndIf
\While{$b-a>1$} \Comment{$\mathrm{pred}(a)$ 거짓, $\mathrm{pred}(b)$ 참}
  \State $m\gets\lfloor a+(b-a)/2\rfloor$; \textbf{if} $m\le a$ \textbf{or} $m\ge b$ \textbf{then} \textbf{break} \Comment{$2^{53}$ 너머}
  \State \textbf{if} $\mathrm{pred}(m)$ \textbf{then} $b\gets m$ \textbf{else} $a\gets m$
\EndWhile
\State \Return $b$
\end{algorithmic}
\caption{갤로핑과 이분으로 찾는 이산 분위수 (stardust \code{csrc/stats/discrete.hpp}의 \code{smallest})}
\label{alg:discrete-icdf}
\end{algorithm}

구현은 어떤 술어를 쓰느냐에 달려 있다. $p\le1/2$이면 $F(k)\ge p$를 쓰지만 $p>1/2$이면 $\bar F(k)\le1-p$를 쓴다. $p\ge1/2$인 double에서 $1-p$는 정확하므로(Sterbenz 보조정리, 11장) 작은 쪽 확률을 그 자체의 상대 정확도로 비교하게 된다. $p<2^{-1000}$이면 확률이 비정규 수에 가까워 상대 정확도를 잃으므로 $\log F(k)\ge\log p$로 비교한다. \code{iccdf}와 \code{ilogccdf}에서는 위 꼬리의 로그가 $-745$ 근처까지 내려가는 분포라면(Zipf처럼 로그 꼬리를 double-double로 줄 수 있는 분포) 그 로그를 double-double로 비교한다. 큰 분위수는 로그의 한 ULP에도 많이 움직이기 때문이다. 이분은 $2^{53}$보다 큰 정수에서 멈춘다. 그 위에서는 이웃한 double 사이에 다른 double이 없어 $b$가 이미 반올림된 답이다. 위로 갤로핑하다 $2^{1023}$을 넘으면 가장 큰 double에서 술어를 한 번 보고 그래도 거짓이면 답을 $\infty$로 둔다. 꼬리가 아주 무거운 Zipf 분포는 이런 분위수를 실제로 갖는다. 갤로핑과 이분은 각각 $1100$번으로 묶여 있어 술어 계산은 많아야 $2\cdot1100+2$번이고 어떤 입력에서도 정해진 시간 안에 끝난다. 분포 함수를 계산할 수 없는 점(꼬리 합이 한도를 넘는 경우)에서는 술어가 계산 없이 거짓을 돌려주므로 탐색은 산술만으로 끝까지 가고 결과는 NaN이 된다.

시작 추측은 탐색의 속도만 정하고 답에는 영향을 주지 않는다. Poisson, 이항, 음이항, 초기하, 베타이항은 표준편차가 $8$ 이상이면 정규 분위수 $z=\Phi^{-1}(p)$에 왜도 $\gamma_1$의 보정을 더한 Cornish--Fisher 근사 $\mu+\sigma(z+\gamma_1(z^2-1)/6)$에서 시작하고 그보다 좁으면 평균 근처에서 시작한다. 기하와 이산 균등처럼 분위수가 닫힌 꼴인 분포는 그 값에서 시작한다. \cref{alg:discrete-icdf}를 그대로 옮겨 Poisson에서 무작위 $p$ $2000$개의 술어 계산 수를 세면 \cref{tab:search-evals}와 같다. Cornish--Fisher에서 시작하면 거의 언제나 두 번으로 끝난다. 평균에서 시작하면 거리 $\ell$이 표준편차 정도여서 계산 수가 $\log_2\sigma$에 비례해 늘어난다.

\begin{table}[htbp]
\begin{sdtable}{colspec={cccccc}, row{1}={font=\small\bfseries}}
\toprule
평균 $\mu$ & $10$ & $40$ & $63.9$ & $64.1$ & $10^4$ \\
\midrule
평균에서 시작 & $4.3$ & $5.7$ & $6.1$ & $6.1$ & $12.5$ \\
Cornish--Fisher에서 시작 & $2.0$ & $2.0$ & $2.0$ & $2.0$ & $2.0$ \\
\bottomrule
\end{sdtable}
\caption{Poisson 분위수 탐색의 술어 계산 수(무작위 $p$ $2000$개의 평균; \cref{alg:discrete-icdf}를 옮긴 코드에 따로 빌드한 stardust의 \code{cdf}·\code{ccdf}를 술어로, \code{figures/09/search\_evals.py}). stardust는 $\sigma\ge8$, 곧 $\mu\ge64$에서 Cornish--Fisher를 쓴다.}
\label{tab:search-evals}
\end{table}

같은 매개변수로 많은 점을 부르는 일이 흔하므로(자료의 분위수, 같은 분포의 표집) stardust는 정수점의 메서드 값을 스레드마다 $1024$칸의 직접 사상 표(\code{KMemo})에 기억해 둔다. 표는 메서드의 정의가 준 값만 돌려주므로 결과의 비트는 기억이 있든 없든 같다. 매개변수가 바뀌면 표를 비운다. 처음 $128$개 원소에서 놓치는 비율이 $75$퍼센트를 넘으면 기억을 끄고 바로 계산한다.

\subsection{Settling at an Atom}

탐색의 비교는 정확한 실수 비교이므로 답은 계산한 $F(k)$의 끝자리까지 정확한 $F^{-}(p)$이다. 탐색의 술어가 사용자가 쓰는 \code{cdf} 메서드 그 자체이면 둘이 같은 double을 비교하므로 $p=$\code{cdf}$(k-1)\le1/2$에서 탐색만으로 $k-1$이 나온다. 술어가 사용자의 함수와 다르면 사정이 달라진다. $p>1/2$이면 탐색은 \code{ccdf}를 $1-p$와 비교한다. 기하 분포처럼 탐색이 닫힌 꼴의 빠른 경로를 쓰는 분포도 있고 로그판은 로그 메서드를 쓴다. 그래서 stardust는 탐색 뒤에 \code{settle}이라는 한 걸음을 더한다. 탐색이 찾은 $k$에 대해 분포 객체 자신이 돌려주는 \code{cdf}$(k-1)$이 이미 $p$ 이상이면 $k-1$을 답으로 한다. 그러면 작은 쪽(확률이 $2^{-1000}$ 이상인 곳)에서는 \code{icdf(cdf(k))==k}가 성립한다(\cref{alg:settle}). 다만 $k-1$이 \cref{def:backward-quantile}의 답이기도 할 때만 옮긴다. $p>1/2$이면 계약의 기준이 $1-p$인데 \code{cdf}$(k-1)$의 반올림은 $1$ 근처에서 $1-p$를 크게 넘어설 수 있다. 평균 $0.0025$인 Poisson에서 $p=1-7.8\times10^{-16}$이면 \code{cdf(4)}가 $p$로 반올림되지만 \code{ccdf(4)}$=8.1\times10^{-16}$이어서 $4$는 계약 밖의 답이다. 이 경우에는 \code{ccdf}$(k-1)\le(1-p)(1+10^{-15})$도 함께 확인한다. 작은 쪽이 $2^{-1000}$ 아래이면 그 double이 상대 정확도를 갖지 않으므로 옮기지 않는다. 비용은 분포 함수 계산 한두 번이다.

\begin{algorithm}[htbp]
\begin{algorithmic}[1]
\Require 목표 $p$, 탐색이 찾은 $k$, 분포의 \code{cdf}, \code{ccdf}
\Ensure 정착한 답
\State \textbf{if} $k-1$이 지지 밖 \textbf{or} $\min(p,1-p)<2^{-1000}$ \textbf{then} \Return $k$
\State \textbf{if} \code{cdf}$(k-1)<p$ \textbf{then} \Return $k$ \Comment{메서드의 double이 이미 $p$에 닿는가}
\State \textbf{if} $p\le1/2$ \textbf{then} \Return $k-1$
\State \textbf{if} \code{ccdf}$(k-1)\le(1-p)(1+10^{-15})$ \textbf{then} \Return $k-1$ \Comment{큰 쪽: 작은 쪽 기준의 계약도}
\State \Return $k$
\end{algorithmic}
\caption{원자에서의 정착 (stardust \code{discrete.hpp}의 \code{settle}, \code{icdf}의 경우)}
\label{alg:settle}
\end{algorithm}

위쪽 절반에서는 \code{icdf(cdf(k))==k}가 보장되지 않는다. 평균 $3$인 Poisson에서 $p=$\code{cdf}$(13)=0.9999965980853868$을 넣으면 $14$가 나온다. 이 $p$는 참값 $F(13)$보다 $1.1\times10^{-17}$ 크므로 정확한 분위수가 바로 $14$이고 계약 안의 답이다. $p$가 $1$에 가까워 반올림하는 사이에 $1-p$가 참값 $\bar F(13)=3.4\times10^{-6}$과 상대 $3\times10^{-12}$만큼 달라졌다. 위쪽 절반에서 짝이 되는 항등식은 \code{iccdf(ccdf(k))==k}이다. 따로 빌드한 stardust로 Poisson 네 개(평균 $0.5$, $3$, $30$, $300$)에서 작은 쪽의 항등식을 모든 $k$에서 시험하고 무작위 $p$ $300$개마다 계약을 정확한 누적 확률로 확인했다(\cref{tab:quantile-tests}).

\begin{table}[htbp]
\begin{sdtable}{colspec={cccc}, row{1}={font=\small\bfseries}}
\toprule
평균 & \code{icdf(cdf(k))==k} ($F\le1/2$) & \code{iccdf(ccdf(k))==k} ($\bar F\le1/2$) & 계약 위반 \\
\midrule
$0.5$ & 해당 없음 & $28/28$ & $0/300$ \\
$3$ & $3/3$ & $40/40$ & $0/300$ \\
$30$ & $30/30$ & $85/85$ & $0/300$ \\
$300$ & $300/300$ & $227/227$ & $0/300$ \\
\bottomrule
\end{sdtable}
\caption{따로 빌드한 stardust의 Poisson 분위수 시험. 항등식은 작은 쪽의 확률이 서로 다른 double인 모든 $k$에서, 계약은 무작위 $p$에서 \cref{def:backward-quantile}을 정확한 누적 확률로 확인했다(\code{figures/09/identities.py}). 무작위 $p$는 원자 근처를 거의 찌르지 않으므로 평균 $0.0025$, $0.5$, $3$, $30$, $300$에서 원자 $F(k)$와 $1-\bar F(k)$의 $\pm1$--$3$ ULP $12386$점도 시험했고 위반은 없었다(\code{contract\_atoms.py}).}
\label{tab:quantile-tests}
\end{table}

연속 분포에서는 근을 찾는다. 정규 분포의 \code{icdf}는 먼저 \code{ndtri}(Wichura의 AS241 유리 근사, 85장)로 표준 분위수 $z$를 구하고 대부분 $\mathrm{fma}(\sigma,z,\mu)$로 바로 끝낸다(\cref{alg:normal-icdf}). \code{ndtri}의 상대오차는 몇 ULP이고 $x=\mu+\sigma z$에서는 그 오차가 $|\sigma z/x|$배가 된다. 그러니 이 비가 $4$ 이하이고 $|z|\le30$이면 걸음이 필요 없다. 따로 빌드한 stardust에서 $N(0,1)$의 \code{icdf}는 $|z|\le30$인 $36053$점 전부에서 \code{special.ndtri}와 비트가 같았다. 걸음은 $\mu+\sigma z$가 소거를 겪을 때와 아주 깊은 꼬리, 그리고 로그정규처럼 $e^{\sigma z}$에서 상대오차가 $|\sigma z|$배로 커지는 때에 쓴다. 그때는 작은 쪽 $q$의 로그에 Newton 한 걸음을 하고 $z$를 double-double로 들고 있다가 $\mu+\sigma z$를 double-double로 만든 뒤 한 번 반올림한다. 걸음 뒤의 오차는 대략 ULP$/|z|$이다. 감마와 베타는 Boost의 시작점과 Halley·Newton 걸음을 쓴다(107장).

\begin{algorithm}[htbp]
\begin{algorithmic}[1]
\Require $p\in(0,1)$, $\mu$, $\sigma$
\Ensure $x=\mu+\sigma\Phi^{-1}(p)$
\State $z_0\gets$ \code{ndtri}$(p)$ \Comment{AS241}
\State \textbf{if} $|\sigma z_0|\le4|\mu+\sigma z_0|$ \textbf{and} $|z_0|\le30$ \textbf{then} \Return $\mathrm{fma}(\sigma,z_0,\mu)$ \Comment{대부분 여기서 끝}
\State $q\gets\min(p,1-p)$; $z\gets$ $q\ge2^{-1000}$ \textbf{?} \code{ndtri}$(q)$ \textbf{:} \code{ndtri\_exp}$(\log q)$
\State $\ell\gets\log\Phi(z)$ ($z\le-1$이면 $-z^2/2$를 TwoProd로, $\log(\mathrm{erfcx}(-z/\sqrt2)/2)$를 더해 double-double)
\State 기울기 $\phi(z)/\Phi(z)$를 erfcx로(지수 둘을 손으로 지운다)
\State $z\gets z+(\log q-\ell)/(\phi/\Phi)$를 double-double로 \Comment{$\log q<-2^{1000}$이면 생략}
\State $p>1/2$이면 $z\gets-z$
\State \Return $\mu+\sigma z$를 double-double로 만들어 한 번 반올림 \Comment{넘침은 반씩 나누어}
\end{algorithmic}
\caption{정규 분포의 분위수 (stardust \code{csrc/stats/dist/normal.hpp})}
\label{alg:normal-icdf}
\end{algorithm}

\section{Sampling by Inversion}
\label{sec:inversion}

\subsection{Uniform Variates from Random Bits}

역변환 표집에는 $(0,1)$의 균등 난수가 필요하지만 컴퓨터는 $64$비트 정수를 만든다. 정수 $m$을 $53$비트로 줄이는 방법은 여럿이고 그 선택이 꼬리의 정확도를 정한다.

\begin{proposition}
\label{prop:uniform53}
$m$이 $\{0,\dots,2^{53}-1\}$에서 균등하고 $u=(m+1/2)\,2^{-53}=(2m+1)\,2^{-54}$을 실수로 보자.
\begin{enumerate}
\item $u$는 $(0,1)$에 있고 $u$와 $1-u$는 같은 분포를 따른다.
\item 작은 쪽 $s=\min(u,1-u)$는 binary64로 정확히 표현되고 가장 작은 값은 $2^{-54}$이다. 반면 $u>1/2$인 $u$는 표현되지 않는다.
\item 모든 $x\in[0,1]$에서 $|P(u\le x)-x|\le2^{-54}$이다.
\end{enumerate}
\end{proposition}

\begin{proof}
$m\mapsto2^{53}-1-m$이 $u\mapsto1-u$를 주는 전단사이므로 대칭이다. 작은 쪽은 $s=(2m'+1)2^{-54}$($m'<2^{52}$)이고 $2m'+1<2^{53}$은 $53$비트 가수에 들어가는 홀수이므로 $s$는 정확히 표현된다. $[1/2,1)$에서 binary64의 간격은 $2^{-53}$이므로 $2^{-54}$의 홀수배는 그 구간에서 표현되지 않는다. 셋째 주장: $P(u\le x)=N/2^{53}$이고 $N=\#\{m:m\le x2^{53}-1/2\}$는 $x2^{53}+1/2$의 바닥을 $[0,2^{53}]$로 자른 값이므로 $|N-x2^{53}|\le1/2$이다.
\end{proof}

흔히 쓰는 $m\,2^{-53}$은 $[0,1)$에 있어 $0$을 낼 수 있고 대칭이 아니다. 역변환에서 $u=0$은 $F^{-}(0)=-\infty$로 간다. 또 가장 작은 양수가 $2^{-53}$이어서 아래 꼬리는 $2^{-53}$보다 깊이 닿지 않는다. \cref{prop:uniform53}의 (ii)에 비추어 보면 구현에서는 $u$를 먼저 double로 만들고 위쪽 반에서 $1-u$를 빼서는 안 된다. 어느 쪽인지와 그 쪽의 작은 값 $\min(u,1-u)$를 따로 만들어야 한다(\cref{alg:by-inverse}). stardust의 표집기는 지금 $u$를 \code{(double(m)+0.5)*0x1p-53}으로 계산한다(균등 분포와 ziggurat의 쐐기 시험은 $[0,1)$의 $m\,2^{-53}$을 쓴다). $m\ge2^{52}$이면 $m+1/2$가 $53$비트에 들지 않아 짝수 쪽으로 반올림되므로 위쪽 반의 $u$는 $2^{-53}$의 짝수배만 나온다. $m=2^{53}-1$이면 $u=1$이 정확히 나와 역생존 함수의 $\bar F^{-}(0)$, 곧 지지의 위 끝(정규 분포라면 $+\infty$)이 확률 $2^{-53}$으로 나온다. 위쪽 꼬리는 $1-u\ge2^{-52}$까지만 닿아 아래쪽보다 얕다. 따로 빌드한 stardust의 표본기 헤더를 그대로 컴파일해 보면 $u=1$에서 Laplace와 로지스틱의 표본은 $+\infty$, Cauchy는 $1.6\times10^{16}$이고 지수 ziggurat의 꼬리는 정확히 경계 $R$의 점질량을 낸다. 정규 ziggurat의 꼬리는 기각되어 해가 없다. \code{random()}은 numpy와 같은 $[0,1)$의 $m\,2^{-53}$을 쓴다.

\subsection{The Inversion Sampler}

\begin{algorithm}[htbp]
\begin{algorithmic}[1]
\Require 난수열의 $64$비트 단어, 분포의 \code{icdf}, \code{iccdf}
\Ensure 분포를 따르는 표본 하나
\State $b\gets$ 단어의 맨 위 비트; $m\gets$ 그다음 $52$비트
\State $s\gets(m+1/2)\,2^{-53}$ \Comment{정확, $s\in(0,1/2)$, \cref{prop:uniform53}}
\If{$b=0$}
  \State \Return \code{icdf}$(s)$ \Comment{아래 반}
\Else
  \State \Return \code{iccdf}$(s)$ \Comment{위 반도 $2^{-54}$까지}
\EndIf
\end{algorithmic}
\caption{역변환 표집(stardust \code{csrc/stats/sample.hpp}의 \code{by\_inverse}를 \cref{prop:uniform53}에 맞게 고친 꼴)}
\label{alg:by-inverse}
\end{algorithm}

\cref{alg:by-inverse}는 \cref{thm:inverse-transform}을 그대로 옮기되 위쪽 반만은 역생존 함수로 계산한다. stardust의 \code{by\_inverse}는 전용 표집기가 없는 분포를 위한 기본 경로인데 지금 단변량 분포 $31$개가 모두 전용 표집기(\code{Custom})가 있어서 실제로 이 경로에 이르는 분포는 없다. 닫힌 꼴 역변환을 쓰는 표집기들(Laplace, Cauchy, 로지스틱 등)이 같은 생각을 분포마다 직접 옮긴다. \code{icdf}$(u)$를 $u>1/2$에서 그대로 부르면 \cref{sec:computing-quantiles}에서 말한 대로 위쪽 꼬리의 정보가 이미 $u$의 반올림에 묻힌다. 작은 쪽 $s$를 정확히 만들어 역생존 함수에 넘기면 위쪽 꼬리도 아래쪽과 같이 $2^{-54}$까지 닿는다. 정규 분포를 역변환으로 뽑는다면 양쪽으로 $|x|\le8.29$까지 나올 수 있다($\Phi^{-1}(2^{-54})=-8.2924$). stardust의 정규 표집기는 ziggurat이고 그 꼬리는 $R-\log(u)/R$ 꼴이라 이론상 약 $14$까지 나간다. 분위수의 후방 오차 계약까지 더하면 표본의 분포 함수는 모든 $x$에서 대략 $|P(X\le x)-F(x)|\le2^{-54}+10^{-14}\min(F(x),1-F(x))$를 만족한다. 가운데에서는 계약의 흔들림이 $u$의 해상도보다 크지만 어디서나 $10^{-14}$의 상대 수준이다.

역변환은 일반적이지만 분위수 계산이 비싼 분포에서는 느리다. stardust는 분위수가 닫힌 꼴인 분포(Laplace, Cauchy, 로지스틱, 지수에서 만드는 Gumbel·Weibull·Pareto·Rayleigh)와 정수 분포의 작은 매개변수(Poisson $\mu<10$, 이항 $np<10$, 최빈값에서 바깥으로 도는 초기하)에서 역변환을 쓴다. 정규·지수에는 ziggurat, 감마에는 Marsaglia--Tsang의 기각법, 큰 Poisson과 이항에는 변환 기각법(PTRS, BTRS)을 쓴다(\cref{tab:samplers}; 105장). 기각법의 정확성도 \cref{thm:inverse-transform}처럼 측도의 등식으로 보이며 그 증명은 105장으로 미룬다.

\begin{table}[htbp]
\begin{sdtable}{colspec={X[2,l]X[3,l]}, row{1}={font=\small\bfseries}, cells={font=\small}}
\toprule
분포 & 방법 \\
\midrule
균등, Laplace, Cauchy, 로지스틱 & 닫힌 꼴의 역변환 \\
Gumbel, GEV, Weibull, Pareto, Rayleigh & 지수 표본의 닫힌 꼴 변환 \\
정규, 로그정규, 지수 & ziggurat($256$층) \\
감마 & Marsaglia--Tsang \\
Poisson, 이항 & 작은 평균에서 역변환, 큰 평균에서 PTRS·BTRS(이항은 $\min(p,1-p)$로 뒤집어) \\
초기하 & 최빈값에서 바깥으로 도는 역변환 \\
기하 & $\lceil\log U/\log(1-p)\rceil$ \\
균등 & $[0,1)$의 $U$로 $a+(b-a)U$($b$로 반올림되면 $a$) \\
Bernoulli, 이산 균등 & Lemire의 유계 정수(2장) \\
카이제곱, $t$, $F$, 역감마 & 감마 표본에서 \\
베타 & $a,b<1$이면 J\"ohnk, 아니면 감마의 비 \\
비중심 카이제곱·$t$·$F$ & 혼합(Poisson 무게의 중심 분포) \\
음이항, 베타이항 & 감마--Poisson, 베타 뒤 이항 \\
로그 급수, Zipf & Kemp의 LK, Devroye의 기각법 \\
\bottomrule
\end{sdtable}
\caption{stardust의 분포별 표집 방법(\code{csrc/stats/sample\_dist.hpp}; 단변량 $31$개가 모두 전용 표집기를 가져 \cref{alg:by-inverse}의 기본 경로에는 이르지 않는다. 105장에서 자세히 다룬다).}
\label{tab:samplers}
\end{table}

\section{Expectation and Moments}
\label{sec:expectation}

\subsection{Integrals against Distributions}

확률변수 $X$의 \term{기댓값}{expectation}은 적분 $E[X]=\int_\Omega X\,dP$이다($E|X|<\infty$일 때). 이 적분은 다음 공식으로 분포 위에서 계산한다.

\begin{proposition}
\label{prop:lotus}
Borel 가측 함수 $g$에 대해 $E|g(X)|<\infty$이면 $E[g(X)]=\int_{\mathbb R}g\,dP_X$이다. 밀도 $f$가 있으면 $\int g(x)f(x)\,dx$이고 정수에 실린 분포이면 $\sum_kg(k)P(X=k)$이다.
\end{proposition}

\begin{proof}
$g=\mathbf 1_B$이면 양변이 모두 $P(X\in B)$이다. 선형성으로 단순 함수에서 성립하고 음이 아닌 $g$는 단순 함수의 증가열의 극한이므로 단조 수렴 정리에서 성립한다. 일반적인 $g$는 $g^+-g^-$로 나눈다.
\end{proof}

음이 아닌 확률변수의 기댓값은 꼬리 확률의 적분이기도 하다. $X\ge0$이면 Tonelli 정리로 $E[X]=E\int_0^\infty\mathbf 1\{X>t\}\,dt=\int_0^\infty P(X>t)\,dt$이고 같은 방법으로 $E[X^k]=\int_0^\infty kt^{k-1}P(X>t)\,dt$이다. 정수 값이면 $E[X]=\sum_{k\ge1}P(X\ge k)$이다. 이 \term{꼬리 적분 공식}{tail integral formula}은 분위수 함수로 쓰면 $E[X]=\int_0^1F^{-}(p)\,dp$가 된다. \cref{thm:inverse-transform}에서 $X$와 $F^{-}(U)$가 같은 분포이기 때문이다.

꼬리 확률을 기댓값으로 묶는 부등식이 셋 있다. $X\ge0$이면 $\mathbf 1\{X\ge a\}\le X/a$의 기댓값을 취해 \term{Markov 부등식}{Markov's inequality} $P(X\ge a)\le E[X]/a$를 얻고 이것을 $(X-\mu)^2$에 쓰면 \term{Chebyshev 부등식}{Chebyshev's inequality} $P(|X-\mu|\ge a)\le\operatorname{Var}X/a^2$이다. 볼록 함수 $\varphi$에 대해서는 \term{Jensen 부등식}{Jensen's inequality} $\varphi(E[X])\le E[\varphi(X)]$가 성립한다. 증명은 이렇다. 8장의 \cref{thm:convex-characterization}에서 $\varphi$는 점 $E[X]$의 지지선 $\ell$ 위에 있으므로 $E[\varphi(X)]\ge E[\ell(X)]=\ell(E[X])=\varphi(E[X])$이다(미분 불가능한 점에서는 부분기울기를 쓴다). Jensen에서 $p\le q$이면 $\|X\|_p\le\|X\|_q$이므로($\varphi(t)=t^{q/p}$) 높은 모멘트가 유한하면 낮은 모멘트도 유한하다. Markov와 Chebyshev는 분포를 모를 때의 한계라서 대개 느슨하다(\cref{fig:tailbounds}). 정규 꼬리는 지수의 제곱으로 줄지만 Chebyshev는 $1/a^2$에 머문다. 그래도 뒤에 나올 큰 수의 법칙처럼 분포를 모르는 채로 수렴을 증명하는 데에는 두 부등식으로 충분하다. 적률생성함수가 있으면 Markov 부등식을 $e^{tX}$에 써서 훨씬 날카로운 한계를 얻는다.

\begin{proposition}[Chernoff]
\label{prop:chernoff}
$K(t)=\log E[e^{tX}]$가 $t\ge0$의 어떤 범위에서 유한하면 $P(X\ge a)\le\exp(-I(a))$, $I(a)=\sup_{t\ge0}(ta-K(t))$이다. 독립이고 같은 분포의 합에서는 $P(S_n\ge na)\le e^{-nI(a)}$이다.
\end{proposition}

\begin{proof}
$t\ge0$이면 Markov에서 $P(X\ge a)\le e^{-ta}E[e^{tX}]=e^{-(ta-K(t))}$이고 $t$에 대해 하한을 취한다. 합에서는 $K_{S_n}=nK$이다.
\end{proof}

$I$는 8장의 켤레 함수 $K^\ast$이다(\cref{sec:subgradients}). 표준 정규에서 $K(t)=t^2/2$이므로 $I(a)=a^2/2$이고 \cref{fig:tailbounds}의 Chernoff 곡선 $2e^{-a^2/2}$이 참값과 같은 지수 속도로 줄어든다. Cram\'er의 대편차 정리에 따르면 이 지수는 날카롭다. 곧 $n^{-1}\log P(S_n\ge na)\to-I(a)$이다. 유계 변수에서는 분포를 몰라도 같은 꼴의 한계가 있다. $X\in[a,b]$, $E[X]=0$이면 $E[e^{tX}]\le e^{t^2(b-a)^2/8}$이다(Hoeffding의 보조정리). 볼록성에서 $e^{tx}\le((b-x)e^{ta}+(x-a)e^{tb})/(b-a)$이고 기댓값을 취한 뒤 그 로그를 $h=t(b-a)$의 함수로 보면 둘째 도함수가 $1/4$ 이하이므로 Taylor 정리에서 $h^2/8$ 이하이다. 이 한계를 Chernoff에 넣으면 \term{Hoeffding 부등식}{Hoeffding's inequality} $P(S_n-E[S_n]\ge s)\le\exp(-2s^2/\sum_i(b_i-a_i)^2)$이 나온다. \cref{sec:ecdf}의 DKW 부등식이 같은 꼴이다.

\begin{figure}[htbp]
\centering
\begin{tikzpicture}
\begin{semilogyaxis}[sdplot,xlabel={$a$},ylabel={꼬리 확률},xmin=0,xmax=8,ymin=1e-15,ymax=1.5,width=0.68\textwidth,height=0.42\textwidth]
\addplot[black,thick] table[x=a,y=normal]{figures/09/tailbounds.dat};
\addlegendentry{$P(|Z|\ge a)$, 정규}
\addplot[black,dashed] table[x=a,y=cheb]{figures/09/tailbounds.dat};
\addlegendentry{Chebyshev $1/a^2$}
\addplot[black!55,thick] table[x=a,y=expo]{figures/09/tailbounds.dat};
\addlegendentry{$P(X\ge a)$, 지수}
\addplot[black!55,dotted,thick] table[x=a,y=markov]{figures/09/tailbounds.dat};
\addlegendentry{Markov $1/a$}
\addplot[black!35,thick,domain=0.1:8,samples=80] {min(1,2*exp(-x^2/2))};
\addlegendentry{Chernoff $2e^{-a^2/2}$}
\end{semilogyaxis}
\end{tikzpicture}
\caption{분포를 모를 때의 한계와 실제 꼬리(표준 정규의 양쪽 꼬리와 Chebyshev, Chernoff; 평균 $1$인 지수와 Markov).}
\label{fig:tailbounds}
\end{figure}

\subsection{Moments, Cumulants, and Generating Functions}

$k$차 \term{모멘트}{moment} $E[X^k]$와 평균 둘레의 \term{중심 모멘트}{central moment} $\mu_k=E[(X-\mu)^k]$는 분포의 모양을 요약한다. $\mu_2=\sigma^2$이 분산이고 $\gamma_1=\mu_3/\sigma^3$을 \term{왜도}{skewness}, $\gamma_2=\mu_4/\sigma^4-3$을 \term{초과 첨도}{excess kurtosis}라 한다. 정규 분포에서는 둘 다 $0$이고 지수 분포에서는 $\gamma_1=2$, $\gamma_2=6$이다.

\term{적률생성함수}{moment generating function} $M(t)=E[e^{tX}]$가 $0$의 근방에서 유한하면 $M(t)=\sum_kE[X^k]t^k/k!$로 모멘트가 모두 나오고 $K(t)=\log M(t)=\sum_k\kappa_kt^k/k!$의 계수를 \term{누적량}{cumulant}이라 한다. $\kappa_1=\mu$, $\kappa_2=\sigma^2$, $\kappa_3=\mu_3$, $\kappa_4=\mu_4-3\sigma^4$이므로 $\gamma_1=\kappa_3/\kappa_2^{3/2}$, $\gamma_2=\kappa_4/\kappa_2^2$이다. $X$와 $Y$가 독립이면 $E[e^{t(X+Y)}]=E[e^{tX}]E[e^{tY}]$이므로 $K_{X+Y}=K_X+K_Y$, 곧 모든 차수의 누적량이 더해진다. 7장의 확률 생성 함수 $E[q^X]$는 $q=e^t$로 놓은 같은 대상이고 \cref{prop:pgf-moments}의 로그 미분이 바로 $\kappa_1$과 $\kappa_2$를 준 계산이었다. 정규 분포 $N(\mu,\sigma^2)$은 $K(t)=\mu t+\sigma^2t^2/2$여서 $3$차 이상의 누적량이 모두 $0$이다.

누적량의 덧셈에서 이미 중심극한정리가 엿보인다. 같은 분포를 따르는 독립 확률변수 $n$개의 합을 표준화한 $S_n^\ast=(S_n-n\mu)/(\sigma\sqrt n)$의 $k$차 누적량은 $k\ge2$에서 $n^{1-k/2}\kappa_k/\sigma^k$이다. $k\ge3$의 누적량은 모두 $0$으로 가고 표준 정규 분포의 누적량만 남는다. 왜도는 $\gamma_1/\sqrt n$, 초과 첨도는 $\gamma_2/n$의 비율로 줄어들며 이 비율이 \cref{sec:clt}의 수렴 속도를 정한다. 표본에서 모멘트를 계산하는 알고리즘(두 번 훑기와 평균의 낮은 부분 보정, 블록 보상 합)은 111장에서 다룬다.

분포 객체의 모멘트는 대개 닫힌 꼴이지만 닫힌 꼴이 소거를 일으키거나 없는 이산 분포에서는 질량 함수를 직접 더한다(\cref{alg:central-moments}). 최빈값에서 시작해 양쪽으로 나가며 중심 모멘트 $m_j=\sum_k(k-\mu)^jp(k)$를 쌓는다. $4$차 항이 비 $r<1$로 줄어들기 시작하면 남은 항을 등비급수 $tr/(1-r)$로 어림하고 이 어림이 $m_4$의 $2^{-60}$배 아래이면 멈춘다. 왜도와 첨도는 $m_3/m_2^{3/2}$, $m_4/m_2^2-3$이다. 초과 첨도의 $-3$은 누적량 $\kappa_4=m_4-3m_2^2$의 소거이고 분포가 정규에 가까울수록 상대 정확도를 잃는다. 그래서 초기하 분포처럼 닫힌 꼴의 괄호가 서로 상쇄되는 경우에는 그 괄호를 $128$비트 정수로 정확히 만든다(108장).

\begin{algorithm}[htbp]
\begin{algorithmic}[1]
\Require 평균 $\mu$, 최빈값, 지지 $[k_{\mathrm{lo}},k_{\mathrm{hi}}]$, $\log p(k)$ (double-double)
\Ensure $m_2$, $m_3$, $m_4$
\State $m_2,m_3,m_4\gets0$; 항 수의 상한 $2^{12}$
\For{$k$를 최빈값에서 위로, 이어 아래로}
  \State $w\gets\exp(\log p(k))$; $d\gets k-\mu$; $m_j\gets m_j+wd^j$ ($j=2,3,4$)
  \State $t\gets wd^4$; $r\gets t/t_{\mathrm{prev}}$
  \State \textbf{if} 최빈값에서 둘 넘게 떨어졌고 $r<1$이며 $tr\le(1-r)m_4\cdot2^{-60}$ \textbf{then} 이 방향을 멈춘다
\EndFor
\State 상한에 닿으면 NaN \Comment{잘린 합은 내지 않는다}
\end{algorithmic}
\caption{이산 분포의 중심 모멘트 (stardust \code{csrc/stats/dist/discrete\_more.hpp}의 \code{central})}
\label{alg:central-moments}
\end{algorithm}

\section{Independence and Product Measures}
\label{sec:independence}

$\sigma$-대수 $\mathcal F_1,\dots,\mathcal F_n$이 \term{독립}{independent}이라는 것은 각각에서 고른 사건 $A_i\in\mathcal F_i$에 대해 $P(A_1\cap\cdots\cap A_n)=\prod_iP(A_i)$인 것이고 확률변수들은 그것들이 생성한 $\sigma$-대수가 독립일 때 독립이다. 무한 개의 족은 유한 부분족이 모두 독립일 때 독립이다. 확인은 생성족에서만 하면 된다.

\begin{proposition}
\label{prop:independence-pi}
$\Omega$를 담는 $\pi$-체계 $\mathcal P_1,\dots,\mathcal P_n$이 서로 독립이면(각각에서 고른 사건들의 곱 공식) 그것들이 생성한 $\sigma$-대수도 독립이다. 특히 확률변수 $X_1,\dots,X_n$은 결합 누적분포함수가 주변 누적분포함수의 곱일 때 그리고 그때에만 독립이다.
\end{proposition}

\begin{proof}
$\Omega\in\mathcal P_i$이므로 곱 공식은 부분족에서도 성립한다. $A_2,\dots,A_n$을 고정하면 곱 공식을 만족하는 $A_1$들의 모임은 $\Omega$를 담는 $\lambda$-체계이고 $\mathcal P_1$을 담으므로 $\pi$--$\lambda$ 정리에서 $\sigma(\mathcal P_1)$을 담는다. 이것을 좌표마다 차례로 되풀이한다. 둘째 주장은 반직선 $(-\infty,x]$의 역상들이 $\pi$-체계라는 데서 나온다.
\end{proof}

독립인 $X$와 $Y$의 결합 분포는 곱측도 $P_X\otimes P_Y$이고 Fubini 정리에서 $E[g(X)h(Y)]=E[g(X)]E[h(Y)]$이다. 합 $X+Y$의 분포는 \term{합성곱}{convolution} $P_X\ast P_Y(B)=\int P_X(B-y)\,dP_Y(y)$이고 밀도가 있으면 $f_{X+Y}(s)=\int f_X(s-y)f_Y(y)\,dy$이다. 7장에서 생성 함수의 곱이 계수의 합성곱이었던 것은 이 공식의 정수판이다. 합성곱은 모양을 매끄럽게 한다. 균등 분포의 밀도는 불연속이지만 둘의 합성곱은 삼각형이고 셋이면 $C^1$인 이차 조각 함수이다. 표준화하면 이 밀도들은 빠르게 정규 밀도에 다가간다(\cref{fig:irwin}). $n$개의 합의 밀도는 Irwin--Hall 공식 $f_n(x)=\sum_j(-1)^j\binom nj(x-j)_+^{n-1}/(n-1)!$이다.

\begin{figure}[htbp]
\centering
\begin{tikzpicture}
\begin{axis}[sdplot,xlabel={표준화한 합 $z$},ylabel={밀도},xmin=-3,xmax=3,ymin=0,ymax=0.5,width=0.68\textwidth,height=0.42\textwidth]
\addplot[black!45,const plot] table[x=z,y=n1]{figures/09/irwin.dat};
\addlegendentry{$n=1$}
\addplot[black,dashed] table[x=z,y=n2]{figures/09/irwin.dat};
\addlegendentry{$n=2$}
\addplot[black!70,thick] table[x=z,y=n3]{figures/09/irwin.dat};
\addlegendentry{$n=3$}
\addplot[black,dotted,thick,domain=-3:3,samples=100] {exp(-x^2/2)/sqrt(2*pi)};
\addlegendentry{정규}
\end{axis}
\end{tikzpicture}
\caption{균등 분포 $n$개의 합을 표준화한 밀도(Irwin--Hall; \code{figures/09/extras.py}). $n=3$이면 이미 정규 밀도와 눈으로 가려지지 않는다.}
\label{fig:irwin}
\end{figure}

난수 생성에서는 독립을 가정으로 받아들이지 않고 설계 목표로 삼는다. stardust의 \code{sd.random}은 카운터 기반 생성기 Philox4x64-10\cite{salmon2011random123}을 쓴다. $128$비트 열쇠로 정해지는 $256$비트 계수기의 전단사 $\Phi_k$이므로 상태가 없고 어느 계수기의 출력이든 바로 계산된다. 원소 $e$(열의 위치와 호출 안의 번호의 합)는 \cref{alg:philox-words}처럼 계수기를 쓴다.

\begin{algorithm}[htbp]
\begin{algorithmic}[1]
\Require 열쇠 $k$, 원소 번호 $e$, 여분 번호 $j\ge1$
\Ensure 원소 $e$의 첫 단어와 $j$번째 여분 블록
\Function{First}{$k,e$} \Comment{네 원소가 한 블록을 나눈다}
  \State \Return $\Phi_k(\lfloor e/4\rfloor,0,0,0)$의 $(e\bmod4)$번째 $64$비트 단어
\EndFunction
\Function{Extra}{$k,e,j$} \Comment{기각법의 다시 뽑기}
  \State \Return $\Phi_k(e,j,0,0)$의 네 단어
\EndFunction
\State 첫 단어들은 $16$원소씩 네 블록을 함께 계산한다 \Comment{지연을 숨긴다}
\end{algorithmic}
\caption{원소별 난수 단어 (stardust \code{csrc/random/philox.hpp}, \code{csrc/stats/sample.hpp}의 \code{Words})}
\label{alg:philox-words}
\end{algorithm}

두 종류의 계수기는 둘째 성분이 $0$인지 $1$ 이상인지로 구별되므로 어떤 단어도 두 번 쓰이지 않는다. 표본은 열쇠, $e$, 매개변수만의 함수이므로 기각 횟수가 원소마다 달라도 다른 원소의 단어를 건드리지 않고 스레드 수나 호출을 나누는 방식과 상관없이 같은 비트가 나온다. 이런 생성기의 출력이 독립이라는 가설은 증명하지 않고 통계적 시험(BigCrush 등)으로 확인한다. 열쇠는 numpy와 같은 \code{SeedSequence}에서 만들고 \code{spawn}은 \code{spawn\_key}에 번호를 붙여 다른 열쇠, 곧 다른 전단사를 준다. 같은 열에서 표본을 나누어 쓸 때는 위치 구간을 원자적으로 예약해 서로 겹치지 않게 한다(104장).

\section{Conditional Expectation}
\label{sec:conditional-expectation}

\subsection{Definition and Existence}

부분 $\sigma$-대수 $\mathcal G\subseteq\mathcal F$와 적분 가능한 $X$에 대해 다음 둘을 만족하는 확률변수 $Y$를 $X$의 $\mathcal G$에 대한 \term{조건부 기댓값}{conditional expectation} $E[X\mid\mathcal G]$라 한다.
\begin{enumerate}
\item $Y$는 $\mathcal G$-가측이고 적분 가능하다.
\item 모든 $G\in\mathcal G$에서 $E[Y\mathbf 1_G]=E[X\mathbf 1_G]$이다.
\end{enumerate}

\begin{theorem}
\label{thm:condexp-exists}
조건부 기댓값은 존재하고 거의 확실히 유일하다.
\end{theorem}

\begin{proof}
유일성: $Y$와 $Y'$이 둘 다 조건을 만족하면 $G=\{Y>Y'+\epsilon\}\in\mathcal G$에서 $0=E[(Y-Y')\mathbf 1_G]\ge\epsilon P(G)$이므로 $P(G)=0$이고 $\epsilon\downarrow0$, 역할을 바꾸어 $Y=Y'$ a.s.이다. 존재: $X\ge0$이라 해도 된다(양과 음의 부분으로 나눈다). $\nu(G)=E[X\mathbf 1_G]$는 $\mathcal G$ 위의 유한 측도이고 $P(G)=0$이면 $\nu(G)=0$이므로 $P|_{\mathcal G}$에 절대 연속이다. Radon--Nikodym 정리가 주는 $\mathcal G$-가측 도함수 $Y=d\nu/dP|_{\mathcal G}$가 (ii)를 만족한다.
\end{proof}

$\mathcal G=\sigma(Z)$이면 $E[X\mid Z]$로 쓴다. Doob--Dynkin 보조정리에서 $\sigma(Z)$-가측인 확률변수는 어떤 Borel 함수 $h$에 대해 $h(Z)$의 꼴이므로 $E[X\mid Z]=h(Z)$이고 이 $h(z)$를 $E[X\mid Z=z]$로 쓴다. $h$는 $P_Z$-거의 어디서나만 정해진다. 결합 밀도 $f(x,z)$가 있으면 $f_Z(z)>0$인 곳에서 조건부 밀도 $f(x\mid z)=f(x,z)/f_Z(z)$로 $h(z)=\int xf(x\mid z)\,dx$이고 (ii)는 Fubini로 바로 확인된다.

\subsection{Properties and Projection}

\begin{proposition}
\label{prop:condexp-properties}
적분 가능한 $X$, $X'$에 대해 거의 확실히 다음이 성립한다.
\begin{enumerate}
\item $E[aX+bX'\mid\mathcal G]=aE[X\mid\mathcal G]+bE[X'\mid\mathcal G]$이고 $X\le X'$이면 $E[X\mid\mathcal G]\le E[X'\mid\mathcal G]$이다.
\item $\mathcal H\subseteq\mathcal G$이면 $E[E[X\mid\mathcal G]\mid\mathcal H]=E[X\mid\mathcal H]$이다(\term{탑 성질}{tower property}). 특히 $E[E[X\mid\mathcal G]]=E[X]$이다.
\item $Z$가 $\mathcal G$-가측이고 $ZX$가 적분 가능하면 $E[ZX\mid\mathcal G]=ZE[X\mid\mathcal G]$이다.
\item $\sigma(X)$가 $\mathcal G$와 독립이면 $E[X\mid\mathcal G]=E[X]$이다.
\end{enumerate}
\end{proposition}

\begin{proof}
모두 오른쪽이 정의의 (i), (ii)를 만족함을 보이고 유일성을 쓴다. (i)의 선형성은 곧바로 나오고 단조성은 유일성의 증명과 같은 논법이다. (ii): $H\in\mathcal H\subseteq\mathcal G$에서 $E[E[X\mid\mathcal G]\mathbf 1_H]=E[X\mathbf 1_H]$이다. (iii): 먼저 $X\ge0$, $Z\ge0$이라 하자. $Z=\mathbf 1_{G'}$($G'\in\mathcal G$)이면 $E[\mathbf 1_{G'}E[X\mid\mathcal G]\mathbf 1_G]=E[X\mathbf 1_{G'\cap G}]$이므로 성립하고 선형성으로 음이 아닌 단순 함수에서 성립한다. 음이 아닌 단순 함수 $Z_k\uparrow Z$를 잡으면 $E[X\mid\mathcal G]\ge0$이므로 $E[Z_kE[X\mid\mathcal G]\mathbf 1_G]=E[Z_kX\mathbf 1_G]$의 양변에 단조 수렴 정리를 쓸 수 있고 $G=\Omega$로 두면 $ZE[X\mid\mathcal G]$가 적분 가능하다. 일반적인 경우는 $Z=Z^+-Z^-$, $X=X^+-X^-$의 네 곱에 이 결과를 쓴다($|Z^\pm X^\pm|\le|ZX|$). (iv): $E[E[X]\mathbf 1_G]=E[X]P(G)=E[X\mathbf 1_G]$(독립)이다.
\end{proof}

$E[X^2]<\infty$이면 조건부 기댓값에 기하적인 뜻이 있다. $L^2(\Omega,\mathcal F,P)$는 Hilbert 공간이고 $\mathcal G$-가측인 제곱 적분 가능한 확률변수들은 닫힌 부분공간 $L^2(\mathcal G)$를 이룬다(그 자체가 완비이다). Hilbert 공간의 사영 정리(8장 \cref{thm:projection}의 증명은 내적만 써서 그대로 옮겨 간다)가 주는 직교 사영 $\Pi X$는 모든 $W\in L^2(\mathcal G)$에서 $E[(X-\Pi X)W]=0$을 만족한다. $W=\mathbf 1_G$로 두면 정의의 (ii)이므로 $E[X\mid\mathcal G]=\Pi X$이다.

\begin{corollary}
\label{cor:best-predictor}
$E[X^2]<\infty$이면 $E[X\mid\mathcal G]$는 $\mathcal G$-가측인 예측 $W$ 가운데 평균 제곱 오차 $E[(X-W)^2]$가 가장 작은 것이다. 특히 $E[(X-E[X\mid Z])^2]\le E[(X-h(Z))^2]$이 모든 Borel 함수 $h$에서 성립한다.
\end{corollary}

\begin{proof}
사영의 직교에서 $E[(X-W)^2]=E[(X-\Pi X)^2]+E[(\Pi X-W)^2]$이다.
\end{proof}

이 따름정리 덕분에 회귀를 조건부 기댓값의 추정으로 볼 수 있다. 같은 직교 분해에서 \term{조건부 분산 공식}{law of total variance} $\operatorname{Var}X=E[\operatorname{Var}(X\mid\mathcal G)]+\operatorname{Var}(E[X\mid\mathcal G])$이 나온다($X-E[X\mid\mathcal G]$와 $E[X\mid\mathcal G]-E[X]$가 직교한다). 붓스트랩과 혼합 분포의 분산, 비중심 분포의 모멘트가 이 공식으로 계산된다. 반응 $Y$와 설명 변수 $Z$가 결합 정규이면 $E[Y\mid Z]$는 $Z$의 일차식이므로 최소제곱 직선이 곧 조건부 기댓값의 추정이다(\cref{fig:condexp}). 일반적인 분포에서 최소제곱은 $Z$의 일차식들이 만드는 더 작은 부분공간에 내린 사영이고 조건부 기댓값은 그것보다 오차가 작거나 같다. 12권의 회귀 모형은 이 사영을 유한 표본에서 계산한다. stardust의 최소제곱은 열 피벗 QR과 반복 개선으로 표본 공간 $\mathbb R^n$에서 같은 사영을 계산한다(123장).

\begin{figure}[htbp]
\centering
\begin{tikzpicture}
\begin{axis}[sdplot,xlabel={$Z$},ylabel={$Y$},xmin=-3.2,xmax=3.2,ymin=-3.2,ymax=3.2,width=0.62\textwidth,height=0.5\textwidth]
\addplot[black!35,only marks,mark=*,mark size=0.6pt] table[x=x,y=y]{figures/09/condexp_pts.dat};
\addlegendentry{표본}
\addplot[black,thick,domain=-3:3] {0.7*x};
\addlegendentry{$E[Y\mid Z]=0.7Z$}
\addplot[black,only marks,mark=square*,mark size=1.6pt] table[x=x,y=y]{figures/09/condexp_bins.dat};
\addlegendentry{구간 평균}
\end{axis}
\end{tikzpicture}
\caption{상관이 $0.7$인 표준 이변량 정규 표본 $400$개(따로 빌드한 stardust)와 조건부 기댓값. $Z$의 구간마다 $Y$의 평균을 내면 조건부 기댓값의 직선 둘레에 모인다(\code{figures/09/condexp.py}).}
\label{fig:condexp}
\end{figure}

\subsection{Gaussian Vectors and Sequential Conditioning}

평균이 $\mu$이고 공분산이 $\Sigma$인 다변량 정규 벡터는 $\Sigma=LL^{\mathsf T}$(Cholesky, 4장)로 분해하면 독립인 표준 정규 $Z$를 써서 $X=\mu+LZ$로 나타난다. stardust의 \code{MultivariateNormal.sample}이 바로 이렇게 뽑는다. 정규 벡터의 조건부 분포는 다시 정규이고 조건부 공분산은 Schur 보수이다.

\begin{proposition}
\label{prop:gaussian-conditional}
$(X_1,X_2)$가 결합 정규이고 $\Sigma_{11}$이 가역이면 $X_2\mid X_1=x_1$은 평균 $\mu_2+\Sigma_{21}\Sigma_{11}^{-1}(x_1-\mu_1)$, 공분산 $\Sigma_{22}-\Sigma_{21}\Sigma_{11}^{-1}\Sigma_{12}$의 정규 분포이다.
\end{proposition}

\begin{proof}
$W=X_2-\Sigma_{21}\Sigma_{11}^{-1}X_1$로 두면 $\operatorname{Cov}(W,X_1)=\Sigma_{21}-\Sigma_{21}\Sigma_{11}^{-1}\Sigma_{11}=0$이다. $(W,X_1)$은 결합 정규이므로 상관이 없으면 독립이다(결합 특성함수가 곱으로 나뉜다). 따라서 $X_2=W+\Sigma_{21}\Sigma_{11}^{-1}X_1$의 조건부 분포는 $W$의 분포를 $\Sigma_{21}\Sigma_{11}^{-1}x_1$만큼 옮긴 것이고 $W$의 공분산을 계산하면 Schur 보수이다.
\end{proof}

다변량 정규 분포 함수 $P(X\le b)$를 계산하는 Genz의 방법\cite{genz1992}은 이 명제와 탑 성질에서 나온다. $X=LY$에서 $L$이 아래 삼각이므로 $X_i\le b_i$는 $Y_i\le(b_i-\sum_{j<i}L_{ij}Y_j)/L_{ii}$이고 이 한계는 앞의 $Y_j$들에만 의존한다. $Y_1,\dots,Y_{i-1}$을 고정한 조건부 확률을 차례로 곱해 나가면 $P(X\le b)=E\bigl[\prod_ie_i\bigr]$이 된다. 여기서 $e_i=\Phi((b_i-\sum_{j<i}L_{ij}y_j)/L_{ii})$이고 $y_j$는 절단된 정규를 역변환으로 뽑은 값이다. 그러면 이 확률은 $(0,1)^{d-1}$ 위의 매끄러운 함수의 적분이 되고 stardust는 이 적분을 준몬테카를로로 계산한다(\cref{alg:genz}).

\begin{algorithm}[htbp]
\begin{algorithmic}[1]
\Require 공분산의 Cholesky 인수 $L$($d\ge4$, 변수 재배열 포함), 상한 $b$, 허용오차 \code{atol}
\Ensure $P(X\le b)$와 보고 오차
\Function{F}{$w\in(0,1)^{d-1}$} \Comment{순차 조건화}
  \State $e_1\gets\Phi(b_1/L_{11})$; $p\gets e_1$; $y_1\gets\Phi^{-1}(w_1e_1)$
  \For{$i=2,\dots,d$}
    \State $e_i\gets\Phi((b_i-\sum_{j<i}L_{ij}y_j)/L_{ii})$; $p\gets pe_i$; $y_i\gets\Phi^{-1}(w_ie_i)$ ($i<d$)
  \EndFor
  \State \Return $p$
\EndFunction
\Repeat \Comment{수준마다 새 격자와 새 이동}
  \State 소수 $N$, 성분별 구성(Nuyens--Cools)의 생성 벡터 $z$
  \For{이동 $m=1,\dots,12$ (Philox, 고정 열쇠)}
    \State $\hat\mu_m\gets N^{-1}\sum_{i<N}F(|2\,\mathrm{frac}(iz/N+\Delta_m)-1|)$ \Comment{빵굽기 변환}
  \EndFor
  \State 수준의 추정 $=\hat\mu_m$들의 평균, 분산 $=$ 그 평균의 분산; 수준들을 역분산 무게로 합친다
\Until{$3\times$(표준오차)$\le$ \code{atol} 또는 계산 상한}
\end{algorithmic}
\caption{다변량 정규 분포 함수의 준몬테카를로 (stardust \code{csrc/stats/mvcdf.hpp}, Genz)}
\label{alg:genz}
\end{algorithm}

\cref{alg:genz}에서 확률론은 오차를 보고하는 일을 맡는다. 무작위 이동 $\Delta_m$이 균등하면 각 $\hat\mu_m$은 참값의 비편향 추정이다. 이동끼리 독립이므로 열두 추정의 평균은 근사적으로 정규이고 그 표본 분산으로 표준화하면 자유도 $11$의 $t$ 분포를 따른다. 표준오차의 $3$배는 이 분포에서 약 $99$퍼센트 구간이다. 격자 규칙의 오차는 $O(N^{-1})$에 가까워 평범한 몬테카를로의 $O(N^{-1/2})$보다 빠르다. 이동의 난수는 고정 열쇠의 Philox이므로 결과는 입력만의 함수이다. 따로 빌드한 stardust로 등상관 공분산($d=4$--$10$, $\rho=0.1$--$0.9$)의 무작위 문제 $300$개를 풀어 일차원 적분으로 구한 참값과 견주었다(\cref{fig:mvn-cover}). 실제 오차는 $300$개 모두 보고 오차와 \code{atol} 안이었고 비율의 중앙값은 $0.23$이었다. 따로 잰 한 경우($d=4$, $\rho=1/2$, $b=0$, \code{atol}$=10^{-6}$)에서는 비율이 $1.07$로 보고 오차를 조금 넘었는데 $99$퍼센트 구간이라는 뜻과 맞는다.

\begin{figure}[htbp]
\centering
\begin{tikzpicture}
\begin{loglogaxis}[sdplot,xlabel={보고한 오차($3\sigma$)},ylabel={실제 오차},xmin=1e-8,xmax=1e-4,ymin=1e-11,ymax=1e-4,width=0.55\textwidth,height=0.5\textwidth]
\addplot[black,only marks,mark=*,mark size=0.8pt] table[x=reported,y=actual]{figures/09/mvn_cover.dat};
\addlegendentry{무작위 문제 $300$개}
\addplot[black!50,dashed,domain=1e-8:1e-4] {x};
\addlegendentry{실제 $=$ 보고}
\end{loglogaxis}
\end{tikzpicture}
\caption{\code{MultivariateNormal.cdf}의 보고 오차와 실제 오차(따로 빌드한 stardust; 참값은 등상관 공분산의 일차원 적분, \code{figures/09/mvn\_cover.py}). 모든 점이 대각선 아래에 있다.}
\label{fig:mvn-cover}
\end{figure}

\section{Modes of Convergence}
\label{sec:convergence}

확률변수의 열 $X_n$이 $X$로 가는 방식은 넷이 있다.
\begin{enumerate}
\item \term{거의 확실한 수렴}{almost sure convergence}: $P(X_n\to X)=1$.
\item \term{확률 수렴}{convergence in probability}: 모든 $\epsilon>0$에서 $P(|X_n-X|>\epsilon)\to0$.
\item $L^p$ 수렴($p\ge1$): $E|X_n-X|^p\to0$.
\item \term{분포 수렴}{convergence in distribution}: $F$가 연속인 모든 점에서 $F_n(x)\to F(x)$. 같은 확률 공간에 있을 필요가 없다.
\end{enumerate}

\begin{proposition}
\label{prop:convergence-implications}
(i)이나 (iii)이면 (ii)이고 (ii)이면 (iv)이다. 거꾸로 (ii)이면 거의 확실히 수렴하는 부분열이 있고 극한이 상수 $c$이면 (iv)에서 (ii)가 나온다.
\end{proposition}

\begin{proof}
(i)$\Rightarrow$(ii): $\mathbf 1\{|X_n-X|>\epsilon\}\to0$ a.s.이고 지배 수렴 정리를 쓴다. (iii)$\Rightarrow$(ii): Markov 부등식에서 $P(|X_n-X|>\epsilon)\le E|X_n-X|^p/\epsilon^p$이다. (ii)$\Rightarrow$(iv): $x$가 $F$의 연속점이면 $P(X_n\le x)\le P(X\le x+\epsilon)+P(|X_n-X|>\epsilon)$이고 $P(X\le x-\epsilon)\le P(X_n\le x)+P(|X_n-X|>\epsilon)$이므로 $n\to\infty$ 뒤 $\epsilon\to0$에서 $F_n(x)\to F(x)$이다. 부분열: $P(|X_{n_k}-X|>2^{-k})\le2^{-k}$인 $n_k$를 잡으면 \cref{lem:borel-cantelli}에서 $|X_{n_k}-X|>2^{-k}$는 유한 번만 일어난다. 상수 극한: $P(|X_n-c|>\epsilon)\le F_n(c-\epsilon)+1-F_n(c+\epsilon/2)\to0$이다.
\end{proof}

다른 함의는 없다. $[0,1]$ 위의 Lebesgue 측도에서 $X_n=n\mathbf 1_{(0,1/n)}$은 거의 확실히 $0$으로 가지만 $E|X_n|=1$이라 $L^1$으로 가지 않는다. 구간을 $[0,1/2],[1/2,1],[0,1/4],\dots$로 훑는 지시 함수의 열은 확률로 $0$에 가지만 모든 점에서 무한히 자주 $1$이므로 거의 확실히 가지 않는다(\cref{fig:convergence}). 분포 수렴은 다음 두 도구와 함께 통계의 근사 이론을 이룬다. 연속 함수 $g$에 대해 $X_n\Rightarrow X$이면 $g(X_n)\Rightarrow g(X)$이다(\term{연속 사상 정리}{continuous mapping theorem}). $X_n\Rightarrow X$이고 $Y_n\to c$(확률)이면 $X_n+Y_n\Rightarrow X+c$, $X_nY_n\Rightarrow cX$이다(\term{Slutsky 정리}{Slutsky's theorem}). 두 정리는 분포 수렴을 ``유계 연속 함수의 기댓값이 수렴한다''로 바꾸어 쓰는 portmanteau 정리에 기댄다. portmanteau 정리는 이 장의 분위수 함수로 곧바로 증명된다.

\begin{proposition}[Skorokhod 표현]
\label{prop:skorokhod}
$X_n\Rightarrow X$이면 $(0,1)$ 위의 균등 $U$에 대해 $Y_n=F_n^{-}(U)$와 $Y=F^{-}(U)$는 각각 $X_n$, $X$와 같은 분포이고 $Y_n\to Y$가 거의 확실하다. 따라서 유계 연속 함수 $g$에 대해 $E[g(X_n)]\to E[g(X)]$이다.
\end{proposition}

\begin{proof}
분포가 같다는 것은 \cref{thm:inverse-transform}에서 나온다. $F^{-}$이 $u$에서 연속이면 $F_n^{-}(u)\to F^{-}(u)$임을 보이면 된다(단조 함수의 불연속점은 가산 개라 측도가 $0$이다). $x<F^{-}(u)$인 $F$의 연속점 $x$를 잡으면 \cref{lem:quantile-galois}에서 $F(x)<u$이고 $F_n(x)\to F(x)$이므로 큰 $n$에서 $F_n(x)<u$, 곧 $F_n^{-}(u)>x$이다. 따라서 $\liminf F_n^{-}(u)\ge F^{-}(u)$이다. $u'>u$를 잡고 $y>F^{-}(u')$인 연속점 $y$를 잡으면 $F(y)\ge u'>u$이므로 큰 $n$에서 $F_n(y)>u$, 곧 $F_n^{-}(u)\le y$이고 $\limsup F_n^{-}(u)\le F^{-}(u')$이다. $u'\downarrow u$에서 연속성을 쓴다. 마지막 주장은 지배 수렴 정리이다.
\end{proof}

연속 사상 정리와 Slutsky 정리는 거의 확실한 수렴에서 자명하므로 이 표현을 거쳐 곧바로 나온다(Slutsky는 $Y_n$이 상수로 가므로 같은 공간에 둘 수 있다). 기댓값의 수렴에는 꼬리를 고르게 묶는 조건이 하나 더 필요하다. $\lim_{M\to\infty}\sup_nE[|X_n|\mathbf 1\{|X_n|>M\}]=0$이면 $X_n$들이 \term{균등 적분 가능}{uniformly integrable}하다고 한다. $X_n\to X$(확률)이고 균등 적분 가능하면 $E|X_n-X|\to0$이다(Vitali). $|X_n|$을 $M$에서 자른 부분은 지배 수렴으로 가고 남은 꼬리는 균등 적분 가능성이 $M$에 대해 고르게 누르기 때문이다. $\sup_nE|X_n|^{1+\delta}<\infty$이면 균등 적분 가능하므로 실제로는 높은 모멘트의 유계로 확인한다. 이 둘에서 \term{델타 방법}{delta method}이 나온다. $\sqrt n(T_n-\theta)\Rightarrow N(0,\sigma^2)$이고 $g$가 $\theta$에서 미분 가능하면 $\sqrt n(g(T_n)-g(\theta))\Rightarrow N(0,g'(\theta)^2\sigma^2)$이다. $r(t)=(g(t)-g(\theta))/(t-\theta)-g'(\theta)$($t\ne\theta$), $r(\theta)=0$으로 두면 $r$은 $\theta$에서 연속이고 $\sqrt n(g(T_n)-g(\theta))=(g'(\theta)+r(T_n))\sqrt n(T_n-\theta)$이다. $T_n-\theta=n^{-1/2}\cdot\sqrt n(T_n-\theta)\to0$(확률, Slutsky)이므로 $r(T_n)\to0$(확률)이고 다시 Slutsky에서 결론이 나온다. 11권에 나오는 추정량의 표준오차는 대부분 이렇게 계산한다.

\begin{figure}[htbp]
\centering
\begin{tikzpicture}[node distance=10mm and 18mm,every node/.style={font=\small},box/.style={draw,rounded corners=2pt,minimum width=24mm,minimum height=7mm}]
\node[box] (as) {거의 확실한 수렴};
\node[box,below=of as] (pr) {확률 수렴};
\node[box,right=of pr] (lp) {$L^p$ 수렴};
\node[box,below=of pr] (d) {분포 수렴};
\draw[-{Stealth},thick] (as) -- (pr);
\draw[-{Stealth},thick] (lp) -- (pr);
\draw[-{Stealth},thick] (pr) -- (d);
\draw[-{Stealth},dashed] (pr.west) to[bend left=40] node[left,align=right,font=\scriptsize] {부분열} (as.west);
\draw[-{Stealth},dashed] (d.east) to[bend right=40] node[right,align=left,font=\scriptsize] {극한이 상수} (pr.east);
\end{tikzpicture}
\caption{수렴의 양식 사이의 함의(\cref{prop:convergence-implications}). 실선은 언제나, 점선은 조건이 붙을 때 성립한다. 표시가 없는 쌍은 반례가 있다.}
\label{fig:convergence}
\end{figure}

\section{Characteristic Functions}
\label{sec:characteristic-function}

적률생성함수는 꼬리가 무거운 분포에서 $0$ 밖의 모든 $t$에서 무한대일 수 있다. 로그정규 분포는 모든 차수의 모멘트가 있지만 적률생성함수가 $t>0$에서 무한대이고 같은 모멘트를 가진 다른 분포가 있다. 그래서 일반적인 도구로는
\[
\varphi(t)=E[e^{itX}]=E[\cos tX]+iE[\sin tX],\qquad t\in\mathbb R
\]
인 \term{특성함수}{characteristic function}를 쓴다. $|e^{itX}|=1$이므로 모든 분포에 대해 존재하고 $|\varphi|\le1$, $\varphi(0)=1$이며 지배 수렴 정리에서 고르게 연속이다. 특성함수는 Fourier 변환(5장, 80장)에서 지수의 부호만 바꾼 것이다. 성질은 셋이다.
\begin{enumerate}
\item 독립 확률변수의 합의 특성함수는 특성함수의 곱이다(Fubini).
\item 특성함수는 분포를 결정한다. 반전 공식
\[
P(a<X<b)+\frac12P(X\in\{a,b\})=\lim_{T\to\infty}\frac1{2\pi}\int_{-T}^T\frac{e^{-ita}-e^{-itb}}{it}\varphi(t)\,dt
\]
이 분포를 되찾는다. 증명은 $\varphi$의 정의를 넣고 Fubini로 순서를 바꾼 뒤 $\int_{-T}^T\sin(t\theta)/t\,dt\to\pi\operatorname{sign}\theta$를 쓰는 것이다(Billingsley\cite{billingsley1995} 정리 26.2).
\item (Lévy의 연속성 정리) $X_n\Rightarrow X$인 것은 모든 $t$에서 $\varphi_n(t)\to\varphi(t)$인 것과 동치이다. 앞쪽은 $e^{itx}$가 유계 연속이라서이고 뒤쪽은 $\varphi_n$의 수렴에서 분포의 열이 조밀함(tightness)을 이끌어 낸 뒤 극한점이 모두 같은 특성함수를 가지므로 (ii)에서 같다는 순서로 보인다. 조밀함의 핵심은 부등식 $P(|X|>2/\delta)\le\delta^{-1}\int_{-\delta}^\delta(1-\varphi(t))\,dt$이다(Fubini로 오른쪽을 $2E[1-\sin(\delta X)/(\delta X)]$로 쓰고 $|\delta X|>2$에서 괄호가 $1/2$ 이상임을 쓴다). $\varphi_n\to\varphi$이고 $\varphi$가 $0$에서 연속이므로 오른쪽이 $n$에 대해 고르게 작아지고 Helly의 선택 정리로 수렴하는 부분열을 얻는다(Billingsley\cite{billingsley1995} 정리 26.3).
\end{enumerate}

$E|X|^k<\infty$이면 $\varphi$는 $k$번 미분 가능하고 $\varphi^{(j)}(0)=i^jE[X^j]$이다. 특히 $E[X]=0$, $E[X^2]=\sigma^2$이면
\begin{equation}
\label{eq:cf-taylor}
\varphi(t)=1-\frac{\sigma^2t^2}2+o(t^2)\qquad(t\to0)
\end{equation}
이다. 이 전개는 $2$차 모멘트만 있으면 성립한다. $e^{ix}$의 Taylor 나머지 부등식 $|e^{ix}-1-ix+x^2/2|\le\min(|x|^3/6,\,x^2)$을 $x=tX$에 쓰고 나머지를 $t^2$로 나눈 $E\min(|t||X|^3/6,X^2)$에 지배 수렴 정리를 쓴다.

표준 정규 분포의 특성함수는 $e^{-t^2/2}$, 표준 Cauchy 분포의 특성함수는 $e^{-|t|}$이다. 그러면 독립인 Cauchy 확률변수 $n$개의 평균의 특성함수는 $(e^{-|t|/n})^n=e^{-|t|}$로 하나의 Cauchy와 같아서 평균을 아무리 많이 내도 흩어짐이 줄지 않는다(\cref{fig:lln}). 특성함수의 $0$ 근처 모양을 보면 모멘트를 알 수 있다(\cref{fig:cf}). 정규와 균등은 $0$에서 매끄럽고(분산이 있다) Cauchy는 $e^{-|t|}$가 $0$에서 꺾여 평균조차 없음을 보인다. 균등 분포의 $\sin t/t$는 $t\to\infty$에서 $1/t$로만 줄어드는데 밀도가 불연속이라서 이렇게 느리게 감쇠한다.

\begin{figure}[htbp]
\centering
\begin{tikzpicture}
\begin{axis}[sdplot,xlabel={$t$},ylabel={$|\varphi(t)|$},xmin=0,xmax=10,ymin=0,ymax=1.05,width=0.68\textwidth,height=0.42\textwidth]
\addplot[black,thick] table[x=t,y=normal]{figures/09/cf.dat};
\addlegendentry{정규 $e^{-t^2/2}$}
\addplot[black,dashed] table[x=t,y=cauchy]{figures/09/cf.dat};
\addlegendentry{Cauchy $e^{-|t|}$}
\addplot[black!55,thick,y filter/.expression={abs(y)}] table[x=t,y=unif]{figures/09/cf.dat};
\addlegendentry{균등$(-1,1)$ $|\sin t/t|$}
\addplot[black!40,only marks,mark=*,mark size=0.5pt] table[x=t,y=emp]{figures/09/cf.dat};
\addlegendentry{경험(정규 표본 $200$개)}
\end{axis}
\end{tikzpicture}
\caption{특성함수의 절댓값. 점은 따로 빌드한 stardust의 정규 표본 $200$개로 만든 경험 특성함수로 큰 $t$에서 $1/\sqrt{200}\approx0.07$ 수준의 잡음에 머문다(\code{figures/09/extras.py}).}
\label{fig:cf}
\end{figure}

특성함수는 자료에서도 쓴다. 표본의 \term{경험 특성함수}{empirical characteristic function} $\hat\varphi(t)=n^{-1}\sum_je^{itx_j}$는 $t$마다 유계 확률변수의 표본평균이므로 \cref{sec:lln}의 법칙에 따라 $\varphi(t)$로 간다. stardust의 \code{htest.epps\_singleton\_2samp}는 두 표본에서 몇 개의 $t_k$에 대한 경험 특성함수를 비교하는 검정이다(117장, \cref{alg:epps}). 꼬리가 무거워도 모든 항이 유계라서 모멘트에 기대는 검정보다 튼튼하다.

\begin{algorithm}[htbp]
\begin{algorithmic}[1]
\Require 표본 $x$ ($n_x$개), $y$ ($n_y$개), $t=(0.4,0.8)$, $N=n_x+n_y$
\Ensure 통계량 $W$와 p값
\State 합친 표본의 사분위 범위로 $\sigma\gets(P_{75}-P_{25})/2$; $s_k\gets t_k/\sigma$
\State $g(v)\gets(\cos s_kv,\ \sin s_kv)_k\in\mathbb R^{2m}$
\State $\bar g_x$, $C_x\gets$ $g(x_i)$의 평균과 편향 공분산; $\bar g_y$, $C_y$도 같게
\State $S\gets(N/n_x)C_x+(N/n_y)C_y$; $d\gets\bar g_x-\bar g_y$
\State $S$의 고유분해(Jacobi)에서 $|\lambda|>10^{-15}\max|\lambda|$인 것만 남겨 의사역행렬 $S^+$, 랭크 $r$
\State $W\gets N\,d^{\mathsf T}S^+d$; $\max(n_x,n_y)<25$이면 작은 표본 보정을 곱한다
\State \Return $W$, $\chi^2_r$의 위 꼬리 확률
\end{algorithmic}
\caption{Epps--Singleton 검정 (stardust \code{csrc/htest/gof2.cpp})}
\label{alg:epps}
\end{algorithm}

\cref{alg:epps}는 \cref{sec:clt}의 다변량 중심극한정리로 정당화된다. 귀무가설 아래 $\sqrt N d$는 $g$의 공분산으로 정해지는 정규 벡터로 가고 $S$가 그 공분산으로 확률 수렴하므로 Slutsky와 연속 사상 정리에서 $W\Rightarrow\chi^2_r$이다. $g$가 유계이므로 어떤 모멘트 조건도 필요 없다.

\section{Laws of Large Numbers}
\label{sec:lln}

같은 분포를 따르는 독립 확률변수 $X_1,X_2,\dots$의 표본평균을 $\bar X_n=S_n/n$이라 하자.

\begin{theorem}[약한 법칙]
\label{thm:wlln}
$E[X_1^2]<\infty$이면 모든 $\epsilon>0$에서 $P(|\bar X_n-\mu|>\epsilon)\to0$이다.
\end{theorem}

\begin{proof}
독립이므로 $\operatorname{Var}\bar X_n=\sigma^2/n$이고 Chebyshev 부등식에서 $P(|\bar X_n-\mu|>\epsilon)\le\sigma^2/(n\epsilon^2)\to0$이다.
\end{proof}

\begin{theorem}[강한 법칙, 4차 모멘트]
\label{thm:slln4}
$E[X_1^4]<\infty$이면 $\bar X_n\to\mu$가 거의 확실하다.
\end{theorem}

\begin{proof}
$\mu=0$이라 해도 된다. $E[S_n^4]=\sum_{i,j,k,l}E[X_iX_jX_kX_l]$에서 첨자 하나가 홀로 나타나는 항은 독립성과 $E[X_i]=0$에서 $0$이다. 남는 것은 네 첨자가 같은 $n$개 항과 두 쌍으로 나뉘는 $3n(n-1)$개 항이므로 $E[S_n^4]=nE[X^4]+3n(n-1)\sigma^4\le Cn^2$이다. Markov 부등식에서
\[
P(|\bar X_n|>\epsilon)=P(S_n^4>\epsilon^4n^4)\le\frac{Cn^2}{\epsilon^4n^4}=\frac C{\epsilon^4n^2}
\]
이고 이 상계의 합은 유한하다. \cref{lem:borel-cantelli}에서 $|\bar X_n|>\epsilon$인 일은 거의 확실히 유한 번만 일어난다. $\epsilon=1/m$($m=1,2,\dots$)의 가산 개 사건을 모으면 $\bar X_n\to0$이 거의 확실하다.
\end{proof}

강한 법칙은 사실 $E|X_1|<\infty$만으로 성립한다. Etemadi는 이를 쌍마다 독립인 경우까지 넓혀 짧게 증명했다\cite{etemadi1981}. 그 줄거리는 이렇다. $X\ge0$으로 줄이고 $Y_k=X_k\mathbf 1\{X_k\le k\}$로 자른 뒤 $\sum_kP(X_k\ne Y_k)=\sum_kP(X>k)\le E[X]$에서 자른 열과 원래 열이 결국 같다. 기하급수적인 부분열 $k_n=\lfloor\alpha^n\rfloor$에서 자른 합 $T_k=\sum_{j\le k}Y_j$에 Chebyshev를 쓰면 $\sum_nP(|T_{k_n}-E[T_{k_n}]|>\epsilon k_n)\le\epsilon^{-2}\sum_n\operatorname{Var}(T_{k_n})/k_n^2\le C\sum_kE[Y_k^2]/k^2\le4CE[X]$이고(쌍마다 독립이면 분산이 더해진다) Borel--Cantelli에서 부분열 위의 수렴을 얻는다. 단조성으로 부분열 사이를 메우고 $\alpha\downarrow1$로 보낸다(Durrett\cite{durrett2019} 2.4절). 반대로 평균이 없으면 강한 법칙은 반드시 깨진다. $E|X_1|=\infty$이면 $\sum P(|X_n|>n)=\infty$이고 \cref{lem:borel-cantelli}의 둘째 반에서 $|X_n|>n$이 무한히 자주 일어나므로 $\bar X_n$이 수렴할 수 없다. \cref{fig:lln}에서 그 대비를 볼 수 있다. 정규와 지수 표본의 누적 평균은 $1/\sqrt n$의 폭으로 좁아지지만 Cauchy 표본의 누적 평균은 $10^6$개 뒤에도 큰 표본 하나가 나올 때마다 뛴다.

\begin{figure}[htbp]
\centering
\begin{tikzpicture}
\begin{semilogxaxis}[sdplot,xlabel={$n$},ylabel={$\bar X_n-\mu$},xmin=1,xmax=1e6,ymin=-2,ymax=2,width=0.75\textwidth,height=0.42\textwidth]
\addplot[black!45] table[x=n,y=cauchy]{figures/09/lln.dat};
\addlegendentry{Cauchy(중앙값 $0$)}
\addplot[black,thick] table[x=n,y=normal]{figures/09/lln.dat};
\addlegendentry{정규}
\addplot[black,dashed] table[x=n,y=exp]{figures/09/lln.dat};
\addlegendentry{지수}
\addplot[black!60,dotted,domain=1:1e6,samples=60] {2/sqrt(x)};
\addlegendentry{$\pm2/\sqrt n$}
\addplot[black!60,dotted,domain=1:1e6,samples=60,forget plot] {-2/sqrt(x)};
\end{semilogxaxis}
\end{tikzpicture}
\caption{따로 빌드한 stardust의 표본(\code{stats.Normal}, \code{Exponential}, \code{Cauchy}의 \code{sample}, 각 $10^6$개)으로 그린 누적 평균. 정규와 지수는 $\pm2/\sqrt n$의 띠 안으로 모이고 Cauchy는 모이지 않는다(\code{figures/09/lln.py}).}
\label{fig:lln}
\end{figure}

몬테카를로 방법은 이 법칙에 기댄다. 기댓값 $E[g(X)]$를 표본평균으로 어림하면 강한 법칙이 수렴을, 다음 절의 중심극한정리가 그 오차 $\sigma_g/\sqrt n$을 준다. stardust의 재표집 검정(122장)은 그 가장 흔한 예이다. 귀무가설 아래 통계량의 분포를 재표집 $n$번으로 어림하고 관측값 이상으로 극단적인 재표집이 $b$개이면 p값을 $(b+1)/(n+1)$로 보고한다. $b/n$ 대신 이 값을 쓰는 까닭은 다음 성질에 있다.

\begin{proposition}
\label{prop:mc-pvalue}
귀무가설 아래 관측 통계량 $T_0$과 재표집 통계량 $T_1,\dots,T_n$이 교환 가능하면(순서를 바꿔도 결합 분포가 같다) $\hat p=(1+\#\{j\ge1:T_j\ge T_0\})/(n+1)$은 모든 $\alpha\in[0,1]$에서 $P(\hat p\le\alpha)\le\alpha$를 만족한다.
\end{proposition}

\begin{proof}
$T$들과 독립인 균등한 무작위 순서로 동률을 깬 뒤 $T_0$이 위에서 몇째인지를 $R$이라 하자. 교환 가능성에서 $R$은 $\{1,\dots,n+1\}$에서 균등하다. 또 $(n+1)\hat p$는 동률을 모두 $T_0$보다 크게 세므로 $R$ 이상이다. 따라서 $P(\hat p\le\alpha)\le P(R\le\alpha(n+1))=\lfloor\alpha(n+1)\rfloor/(n+1)\le\alpha$이다.
\end{proof}

$b/n$은 관측값보다 극단적인 재표집이 하나도 없으면 $0$이 되어 이 성질을 깬다. stardust는 \cref{alg:resample-p}처럼 계산한다. 반올림 오차가 있어도 관측값과 같은 값을 알아보도록 상대 $100\epsilon$의 여유를 두고 그 안의 값을 극단 쪽으로 센다. \cref{prop:mc-pvalue}가 동률을 극단 쪽으로 세어도 성립하므로 이 여유는 계약을 지킨다. 따로 빌드한 stardust의 \code{permutation\_test}($n=99$)를 귀무가설 아래에서 $400$번 돌려 보니 p값의 경험 분포는 \cref{fig:perm-p}처럼 대각선 위에 있었고 $P(\hat p\le0.05)$는 $0.060$(표준오차 $0.011$)이었다.

\begin{algorithm}[htbp]
\begin{algorithmic}[1]
\Require 귀무 분포의 표본 $T_1,\dots,T_n$, 관측값 $T_0$, 대립가설, 전수 나열인지
\Ensure p값
\State $a\gets$ (전수 나열이면 $0$, 아니면 $1$); $\gamma\gets100\epsilon|T_0|$
\State $\ell\gets\#\{T_j\le T_0+\gamma\}$; $h\gets\#\{T_j\ge T_0-\gamma\}$
\State $p_\ell\gets(\ell+a)/(n+a)$; $p_h\gets(h+a)/(n+a)$
\State \Return 대립가설에 따라 $p_\ell$, $p_h$, 또는 $\min(1,2\min(p_\ell,p_h))$
\end{algorithmic}
\caption{재표집 p값 (stardust \code{csrc/htest/resampling.cpp})}
\label{alg:resample-p}
\end{algorithm}

\begin{figure}[htbp]
\centering
\begin{tikzpicture}
\begin{axis}[sdplot,xlabel={$\alpha$},ylabel={$P(\hat p\le\alpha)$의 경험값},xmin=0,xmax=1,ymin=0,ymax=1,width=0.5\textwidth,height=0.45\textwidth]
\addplot[black,const plot,thick] table[x=p,y=ecdf]{figures/09/perm_p.dat};
\addlegendentry{\code{permutation\_test}}
\addplot[black!50,dashed,domain=0:1] {x};
\addlegendentry{균등}
\end{axis}
\end{tikzpicture}
\caption{귀무가설 아래의 순열 검정 p값($15$개씩 두 정규 표본, 재표집 $99$번, $400$번 반복; 따로 빌드한 stardust, \code{figures/09/perm\_pvalue.py}). p값은 $1/100$의 격자에 있고 경험 분포가 대각선 근처에 머문다.}
\label{fig:perm-p}
\end{figure}

모든 순열을 다 나열하는 정확한 순열 검정에서는 표본평균으로 어림하지 않고 정확한 분포를 얻으므로 $b/n$을 쓴다(이때 $b$는 항등 순열, 곧 관측 배치를 포함해 센다. 따로 빌드한 stardust에서 $(5,6,7)$과 $(1,2,3)$의 정확한 단측 p값은 $20$가지 배치 가운데 하나인 $1/20$이었다). 붓스트랩의 표준오차는 붓스트랩 분포의 표준편차, 곧 큰 수의 법칙으로 수렴하는 표본 표준편차이다.

\section{The Central Limit Theorem}
\label{sec:clt}

\begin{theorem}[Lindeberg--L\'evy]
\label{thm:clt}
$X_i$들이 독립이고 같은 분포를 따르며 $E[X_1]=\mu$, $0<\operatorname{Var}X_1=\sigma^2<\infty$이면 $(S_n-n\mu)/(\sigma\sqrt n)\Rightarrow N(0,1)$이다.
\end{theorem}

\begin{proof}
$\mu=0$, $\sigma=1$이라 해도 된다. 표준화한 합의 특성함수는 $\varphi(t/\sqrt n)^n$이고 \cref{eq:cf-taylor}에서 $\varphi(t/\sqrt n)=1-t^2/(2n)+o(1/n)$이다. 절댓값이 $1$ 이하인 복소수 $a_j$, $b_j$에 대해 $|\prod a_j-\prod b_j|\le\sum|a_j-b_j|$이므로(항을 하나씩 바꾸는 망원합) $n$이 커서 $|1-t^2/(2n)|\le1$이면
\[
\Bigl|\varphi(t/\sqrt n)^n-\Bigl(1-\frac{t^2}{2n}\Bigr)^n\Bigr|\le n\cdot o(1/n)\to0
\]
이다. $(1-t^2/(2n))^n\to e^{-t^2/2}$이므로 특성함수가 표준 정규 분포의 특성함수로 수렴하고 L\'evy의 연속성 정리에서 분포가 수렴한다.
\end{proof}

Lindeberg는 1922년에 분포가 같지 않아도 각 항이 합의 분산에 견주어 무시할 만하면 같은 결론이 성립함을 보였다\cite{lindeberg1922}. 평균 $0$, 분산 $\sigma_i^2$, $s_n^2=\sum_{i\le n}\sigma_i^2$일 때 그 조건은 모든 $\epsilon>0$에서
\[
\frac1{s_n^2}\sum_{i=1}^nE\bigl[X_i^2\mathbf 1\{|X_i|>\epsilon s_n\}\bigr]\to0
\]
이다. 증명은 위와 같은 특성함수의 곱을 항마다 묶는다. $|\varphi_i(t/s_n)-(1-t^2\sigma_i^2/(2s_n^2))|\le E\min(|tX_i|^3/(6s_n^3),\,t^2X_i^2/s_n^2)$에서 $|X_i|\le\epsilon s_n$인 부분은 $\epsilon|t|^3\sigma_i^2/s_n^2$으로, 나머지는 $t^2E[X_i^2\mathbf 1\{|X_i|>\epsilon s_n\}]/s_n^2$으로 누르고 $i$에 대해 더하면 앞은 $\epsilon|t|^3$, 뒤는 Lindeberg 조건에서 $0$으로 간다(Billingsley\cite{billingsley1995} 정리 27.2). 7장의 \cref{prop:rank-clt}에서 쓴 Lyapunov 조건은 이 조건의 충분조건이다. 회귀 계수의 추정량처럼 서로 다른 무게의 합에 정규 근사를 쓸 때 이 형태가 필요하다. 벡터로 넓힐 때는 \term{Cram\'er--Wold 장치}{Cram\'er--Wold device}를 쓴다. $\mathbb R^d$에서 $X_n\Rightarrow X$인 것은 모든 $a\in\mathbb R^d$에서 $a^{\mathsf T}X_n\Rightarrow a^{\mathsf T}X$인 것과 같다. $X$의 특성함수가 $\varphi_X(a)=E[e^{ia^{\mathsf T}X}]=\varphi_{a^{\mathsf T}X}(1)$이므로 L\'evy의 연속성 정리(다변량판)에서 바로 나온다. 그러면 독립이고 같은 분포의 벡터의 합에 대한 다변량 중심극한정리 $\sqrt n(\bar X_n-\mu)\Rightarrow N(0,\Sigma)$는 일차원 정리를 모든 $a^{\mathsf T}X_i$에 쓰면 나오고 다변량 델타 방법도 같은 방식으로 얻는다. \cref{alg:epps}의 $\chi^2$ 극한과 11--12권의 Wald 통계량이 이 다변량 정리에 기댄다.

\subsection{Rate of Convergence}

중심극한정리 자체는 수렴이 얼마나 빠른지 알려 주지 않는다. Berry\cite{berry1941}와 Esseen은 $3$차 절대 모멘트 $\rho=E|X_1-\mu|^3$이 유한하면
\[
\sup_x\bigl|F_n(x)-\Phi(x)\bigr|\le\frac{C\rho}{\sigma^3\sqrt n}
\]
임을 보였다. $F_n$은 표준화한 합의 누적분포함수이고 같은 분포의 경우 $C\le0.4748$이 알려져 있다\cite{shevtsova2011}. 증명의 줄거리는 Esseen의 평활 부등식 $\sup|F-G|\le\pi^{-1}\int_{-T}^T|\varphi_F-\varphi_G|/|t|\,dt+24\sup|G'|/(\pi T)$로 분포의 거리를 특성함수의 거리로 바꾸고 $|t|\le c\sqrt n$에서 $\varphi(t/\sqrt n)^n$과 $e^{-t^2/2}$의 차를 $C\rho|t|^3e^{-t^2/4}/(\sigma^3\sqrt n)$으로 묶는 것이다. 이 $e^{-t^2/4}$ 인자가 있어야 $|t|^{-1}$의 무게로 적분해 $1/\sqrt n$을 얻는다(Petrov\cite{petrov1975} 5장).

이 경계가 얼마나 날카로운지 지수 분포로 확인했다. 평균 $1$인 지수 확률변수 $n$개의 합은 모양 $n$의 감마 분포이므로 $F_n(x)=1-e^{-y}\sum_{k<n}y^k/k!$($y=n+x\sqrt n$)을 십진 산술로 정확히 계산할 수 있다. $\rho=12/e-2\approx2.4146$이고 경계는 $1.146/\sqrt n$이다. 실제 최대 오차에 $\sqrt n$을 곱하면 $n=1,3,10,100,10^4$에서 $0.159$, $0.1335$, $0.133$, $0.1330$, $0.1330$으로 모이고 경계보다 약 $8.6$배 작다(\cref{fig:clt}). 상수 $0.1330$은 \term{Edgeworth 전개}{Edgeworth expansion}의 첫 보정항에서 나온다. $3$차 모멘트가 있고 분포가 격자가 아니면 아래 식이 $o(1/\sqrt n)$까지 맞는다(Esseen\cite{esseen1945}). Cram\'er 조건 $\limsup_{|t|\to\infty}|\varphi(t)|<1$(밀도가 있으면 Riemann--Lebesgue 보조정리로 성립한다)과 $4$차 모멘트가 있으면 $O(1/n)$까지 맞는다. 누적량으로
\[
F_n(x)=\Phi(x)-\frac{\gamma_1}{6\sqrt n}(x^2-1)\phi(x)+O(1/n)
\]
이다. 유도는 누적량에서 시작한다. 표준화한 합의 $\log\varphi_n(t)=-t^2/2+\kappa_3(it)^3/(6\sigma^3\sqrt n)+O(t^4/n)$이므로 $\varphi_n(t)\approx e^{-t^2/2}(1+\gamma_1(it)^3/(6\sqrt n))$이고 $(it)^ke^{-t^2/2}$의 역변환이 $(-D)^k\phi$이므로 $F_n\approx\Phi-(\gamma_1/(6\sqrt n))\Phi'''=\Phi-\gamma_1(x^2-1)\phi/(6\sqrt n)$이다. 나머지는 Esseen의 평활 부등식으로 묶는다(Petrov\cite{petrov1975} 6장). 보정항의 크기는 $x=0$에서 가장 커서 $\gamma_1\phi(0)/6=2\times0.39894/6=0.13298$이다. 최대 오차의 위치도 $n$이 커지면서 $0$으로 다가갔다($n=10$에서 $x=-0.026$, $10^4$에서 $-0.001$). 오차는 주로 왜도에서 온다. 그래서 대칭이면서 격자가 아닌 분포에서는 정규 근사가 $1/n$의 비율로 더 빨리 좋아진다. 동전 던지기 같은 격자 분포는 대칭이어도 계단 때문에 오차가 $1/\sqrt n$에 머물고 연속성 보정이 그 계단을 반으로 나눈다(7장 \cref{sec:normal-approx}).

\begin{figure}[htbp]
\centering
\begin{tikzpicture}
\begin{loglogaxis}[sdplot,xlabel={$n$},ylabel={$\sup_x|F_n(x)-\Phi(x)|$},xmin=0.8,xmax=1.3e4,ymin=8e-4,ymax=2]
\addplot[black,mark=*,mark size=1.2pt] table[x=n,y=err]{figures/09/clt_gamma.dat};
\addlegendentry{정확한 최대 오차}
\addplot[black,dashed] table[x=n,y=be]{figures/09/clt_gamma.dat};
\addlegendentry{Berry--Esseen 경계}
\addplot[black!50,dotted,thick] table[x=n,y=edg]{figures/09/clt_gamma.dat};
\addlegendentry{Edgeworth 첫 항 $0.133/\sqrt n$}
\end{loglogaxis}
\end{tikzpicture}
\caption{지수 분포 $n$개의 합을 표준화한 분포와 표준 정규 분포의 최대 거리(감마 분포 함수를 십진 산술로 정확히 계산, \code{figures/09/clt\_gamma.py}). $n\ge10$에서 정확한 값과 Edgeworth 첫 항이 겹친다.}
\label{fig:clt}
\end{figure}

\begin{aside}[정규 분포와 오차의 법칙]
정규 분포는 de Moivre가 1733년에 이항 분포의 근사로 처음 얻었고 Laplace가 일반적인 합으로 넓혔으며 Gauss가 측정 오차의 분포로 썼다. ``중심극한정리''라는 이름은 확률론에서 중심이 되는 극한 정리라는 뜻으로 P\'olya가 1920년에 붙였다\cite{polya1920}. Lyapunov가 1901년에 특성함수로 처음 엄밀한 일반 증명을 냈고 Lindeberg의 1922년 조건과 L\'evy의 연속성 정리로 오늘날의 꼴이 갖추어졌다.
\end{aside}

\section{The Empirical Distribution Function}
\label{sec:ecdf}

표본 $X_1,\dots,X_n$의 \term{경험 분포 함수}{empirical distribution function} $F_n(x)=n^{-1}\#\{i:X_i\le x\}$는 $x$마다 지시 함수의 표본평균이므로 강한 법칙에서 $F_n(x)\to F(x)$ a.s.이다. 이 수렴이 $x$에 대해 고르게 일어난다는 정리를 통계학의 기본 정리라 부른다.

\begin{theorem}[Glivenko--Cantelli]
\label{thm:glivenko-cantelli}
$D_n=\sup_x|F_n(x)-F(x)|\to0$이 거의 확실하다.
\end{theorem}

\begin{proof}
$\epsilon>0$과 정수 $k>1/\epsilon$을 잡고 $x_j=F^{-}(j/k)$($j=1,\dots,k-1$), $x_0=-\infty$, $x_k=\infty$로 두고 $F(x_0)=F_n(x_0)=0$, $F(x_k-)=F_n(x_k-)=1$로 둔다. 강한 법칙을 유한 개의 점 $x_j$에서 $F_n(x_j)$와 $F_n(x_j-)$에 쓰면 거의 확실히 모든 $j$에서 $F_n(x_j)\to F(x_j)$, $F_n(x_j-)\to F(x_j-)$이다. \cref{lem:quantile-galois}에서 $F(x_{j}-)\le j/k\le F(x_j)$이므로 $F(x_{j+1}-)-F(x_j)\le1/k<\epsilon$이다. $x\in[x_j,x_{j+1})$이면 단조성에서
\[
F_n(x)-F(x)\le F_n(x_{j+1}-)-F(x_j)\le F_n(x_{j+1}-)-F(x_{j+1}-)+\epsilon
\]
이고 같은 방법으로 아래쪽도 묶이므로 $\limsup D_n\le\epsilon$ a.s.이다. $\epsilon=1/m$의 가산 개를 모은다.
\end{proof}

속도에 관해서는 더 정확한 결과가 있다. Dvoretzky, Kiefer, Wolfowitz\cite{dvoretzky1956}가 보이고 Massart\cite{massart1990}가 상수를 최적으로 만든 부등식 $P(D_n>t)\le2e^{-2nt^2}$은 모든 $n$에서 성립한다. 또 $F$가 연속이면 $\sqrt nD_n$은 분포에 의존하지 않고 Kolmogorov 분포 $P(\sqrt nD_n\le t)\to1-2\sum_{k\ge1}(-1)^{k-1}e^{-2k^2t^2}$로 간다\cite{kolmogorov1933ks}. 분포에 의존하지 않는 것은 \cref{thm:inverse-transform}에서 $F(X_i)$가 균등하므로 $D_n$이 균등 표본의 같은 양과 같은 분포를 갖기 때문이다. DKW의 줄거리도 여기서 시작한다. 균등 표본으로 줄이면 한쪽 통계량 $D_n^+$의 정확한 분포가 \cref{alg:ks-stat} 뒤에 적은 Birnbaum--Tingey 합이고 Massart는 그 합의 항들을 묶어 상수 $2$를 얻었다\cite{massart1990}. 상수를 따지지 않는 약한 판은 대칭화와 \cref{prop:chernoff} 뒤의 Hoeffding 부등식으로 증명할 수 있다.

Kolmogorov 극한의 생존 함수 $Q(x)=2\sum_{k\ge1}(-1)^{k-1}e^{-2k^2x^2}$에는 다른 꼴이 하나 더 있다. 5장의 Poisson 합 공식을 Gauss 함수에 쓰면(Jacobi의 세타 변환)
\[
\sum_{k\in\mathbb Z}(-1)^ke^{-2k^2x^2}=\frac{\sqrt{2\pi}}x\sum_{k\ge1}e^{-(2k-1)^2\pi^2/(8x^2)}
\]
이고 왼쪽은 $1-Q(x)$이다. 첫 꼴은 큰 $x$에서, 둘째 꼴은 작은 $x$에서 몇 항이면 수렴한다. stardust의 \code{kolmogorov\_sf}가 $x=1$에서 두 꼴을 나누어 쓴다(\cref{alg:kolmogorov}).

\begin{algorithm}[htbp]
\begin{algorithmic}[1]
\Require $x>0$
\Ensure $Q(x)=P(\sqrt nD_n>x)$의 극한
\If{$x<1$}
  \State $w\gets\pi^2/(8x^2)$; \textbf{if} $w>745$ \textbf{then} \Return $1$
  \State $s\gets\sum_{k\ge1}e^{-(2k-1)^2w}$ (항이 합의 $2^{-60}$배 아래가 되면 멈춘다)
  \State \Return $1-\sqrt{2\pi}\,s/x$
\Else
  \State $x^2=h+\ell$ (TwoProd) \Comment{지수의 큰 인자를 정확히}
  \State \Return $2\sum_{k\ge1}(-1)^{k-1}e^{-2k^2h}(1-2k^2\ell)$ (같은 멈춤)
\EndIf
\end{algorithmic}
\caption{Kolmogorov 분포의 생존 함수 (stardust \code{csrc/htest/ks.hpp}의 \code{kolmogorov\_sf})}
\label{alg:kolmogorov}
\end{algorithm}

큰 $x$에서는 $2x^2$이 커서 $e^{-2k^2x^2}$의 지수에 $x^2$의 반올림 오차가 $2k^2\cdot|x^2|u$만큼 들어가고 그만큼 상대오차가 커진다. $x^2$을 TwoProd로 $h+\ell$로 나누어 $e^{-2k^2h}(1-2k^2\ell)$로 계산하면 이 오차가 사라진다. \cref{alg:lse}의 TwoSum 보정과 같은 생각이다. 극한은 이렇게 얻는다. 균등 경험 과정 $\sqrt n(F_n(u)-u)$가 Brown 다리로 가므로 7장의 반사 원리를 연속판으로 옮겨 그 최댓값의 분포를 구한다(\cref{prop:two-sided-reflection}의 교대합이 같은 꼴이다). \cref{fig:gc}에서 정규 표본의 $D_n$은 $1/\sqrt n$으로 줄고 $\sqrt nD_n$은 Kolmogorov 분포의 범위($0.4$--$1.4$ 정도)에서 흔들린다. 한 열의 앞부분을 본 것이라 이웃한 $n$의 값끼리는 서로 상관되어 있다.

\begin{figure}[htbp]
\centering
\begin{tikzpicture}
\begin{loglogaxis}[sdplot,xlabel={$n$},ylabel={$D_n$},xmin=8,xmax=1.3e6,ymin=3e-4,ymax=1,width=0.47\textwidth,height=0.4\textwidth]
\addplot[black,mark=*,mark size=1.1pt] table[x=n,y=D]{figures/09/gc.dat};
\addplot[black!50,dashed,domain=8:1.3e6,samples=10] {0.87/sqrt(x)};
\end{loglogaxis}
\end{tikzpicture}\hfill
\begin{tikzpicture}
\begin{semilogxaxis}[sdplot,xlabel={$n$},ylabel={$\sqrt nD_n$},xmin=8,xmax=1.3e6,ymin=0,ymax=2,width=0.47\textwidth,height=0.4\textwidth]
\addplot[black,mark=*,mark size=1.1pt] table[x=n,y=sqrtnD]{figures/09/gc.dat};
\addplot[black!50,dashed] coordinates {(8,0.87) (1.3e6,0.87)};
\end{semilogxaxis}
\end{tikzpicture}
\caption{정규 표본(따로 빌드한 stardust, 하나의 열의 앞 $n$개)의 $D_n=\sup|F_n-\Phi|$(왼쪽)와 $\sqrt nD_n$(오른쪽). 점선은 Kolmogorov 분포의 평균 $\sqrt{\pi/2}\,\ln2\approx0.87$을 $\sqrt n$으로 나눈 것이다(\code{figures/09/lln.py}).}
\label{fig:gc}
\end{figure}

순서통계량 자체의 분포도 같은 생각으로 얻는다. $k$번째로 작은 값 $X_{(k)}$가 $x$ 이하인 것은 $x$ 이하의 표본이 $k$개 이상인 것이고 그 개수는 이항 분포 $\mathrm{Bin}(n,F(x))$를 따르므로
\[
P(X_{(k)}\le x)=\sum_{j=k}^n\binom njF(x)^j(1-F(x))^{n-j}
\]
이다. 균등 표본이면 미분해 $X_{(k)}$의 밀도 $n!/((k-1)!(n-k)!)\,u^{k-1}(1-u)^{n-k}$, 곧 베타 분포 $\mathrm{Beta}(k,n-k+1)$을 얻는다. 위 합이 정규화된 불완전 베타 함수와 같다는 항등식(88장)이 바로 이 관계이다. 표본 분위수의 신뢰구간(111장)과 KS의 정확한 분포도 모두 여기에 기댄다.

$D_n$은 순서통계량에서 정확히 계산된다. $F$가 연속이고 정렬한 표본이 $x_{(1)}<\dots<x_{(n)}$이면 $F_n-F$는 각 $x_{(i)}$ 바로 뒤와 바로 앞에서 극값을 가지므로
\[
D_n^+=\max_i\Bigl(\frac in-F(x_{(i)})\Bigr),\qquad D_n^-=\max_i\Bigl(F(x_{(i)})-\frac{i-1}n\Bigr),\qquad D_n=\max(D_n^+,D_n^-)
\]
이다. stardust의 \code{htest.ks\_1samp}(115장)는 정렬 한 번과 $F$의 계산 $n$번, 한 번의 훑기로 두 최댓값과 그 위치를 함께 찾는다(\cref{alg:ks-stat}).

\begin{algorithm}[htbp]
\begin{algorithmic}[1]
\Require 표본 $x_1,\dots,x_n$, 연속 분포 함수 $F$
\Ensure $D_n^+$, $D_n^-$와 그 위치
\State 표본을 정렬하고(NaN은 정책대로) $f_i\gets F(x_{(i)})$
\State $d^+\gets-\infty$; $d^-\gets-\infty$
\For{$i=1,\dots,n$}
  \State \textbf{if} $i/n-f_i>d^+$ \textbf{then} $d^+\gets i/n-f_i$, 위치 $x_{(i)}$를 기억
  \State \textbf{if} $f_i-(i-1)/n>d^-$ \textbf{then} $d^-\gets f_i-(i-1)/n$, 위치 $x_{(i)}$를 기억
\EndFor
\State \Return $d^+$, $d^-$ (양쪽 검정이면 큰 쪽과 그 부호)
\end{algorithmic}
\caption{한 표본 Kolmogorov--Smirnov 통계량 (stardust \code{csrc/htest/gof.cpp}의 \code{ks\_test})}
\label{alg:ks-stat}
\end{algorithm}

위치 $i$의 분수 $i/n$에는 정수 $i$를 double로 나눌 때의 반올림 한 번만 들어가고 $F$의 계산이 상대 $10^{-14}$이므로 $D_n$의 오차는 $10^{-14}$ 남짓이다. 양쪽 검정에서 $D_n^+=D_n^-$이면 $D_n^-$ 쪽을 고른다. p값의 경로는 셋이다. 정확한 경로는 Ruben--Gambino의 닫힌 꼴, $nd^2\ge6.5$이거나 $d\ge1/2$이면 $2\cdot$(한쪽 꼬리), 그 밖에는 Durbin의 행렬 거듭제곱(Marsaglia--Tsang--Wang)을 double-double로 계산하고 일이 상한을 넘으면 Pelz--Good 근사로 간다. 한쪽 꼬리는 Birnbaum--Tingey의 합 $P(D_n^+\ge d)=d\sum_{j\le n(1-d)}\binom nj(1-d-j/n)^{n-j}(d+j/n)^{j-1}$을 double-double 로그로 더한다. 점근 경로는 Kolmogorov 극한이다(\cref{alg:kolmogorov}).

\section{Numerical Computation of Probabilities}
\label{sec:probability-numerics}

수학적으로는 $P(X>x)=1-F(x)$이지만 binary64로 그렇게 계산하면 $F(x)$가 $1$에 가까울 때 소거가 일어난다(11장). 표준 정규 분포에서 $\Phi(5)=0.9999997133484281$의 반올림 오차는 $u$ 정도인데 $1-\Phi(5)$로 빼면 그 오차가 $1/(1-\Phi(5))\approx3.5\times10^6$배로 커져 상대오차 $1.5\times10^{-10}$이 된다. $x=8$이면 $7$퍼센트이고 $x=10$이면 $\Phi(10)$이 정확히 $1$로 반올림되어 $1-\Phi(10)=0$이지만 참값은 $7.62\times10^{-24}$이다. 그래서 수치 라이브러리는 누적분포함수와 따로 \term{생존 함수}{survival function} $\bar F(x)=1-F(x)$를 계산한다.

\begin{algorithm}[htbp]
\begin{algorithmic}[1]
\Require 분포의 작은 쪽을 직접 계산하는 함수(\code{cdf}의 아래 꼬리, \code{ccdf}의 위 꼬리), 점 $x$, 원하는 것(\code{cdf}, \code{ccdf}, 그 로그)
\Ensure 상대 정확한 값
\State $s\gets$ $x$가 있는 쪽의 작은 꼬리 확률(직접 계산) \Comment{$s\le1/2$}
\If{원하는 것이 작은 쪽} \Return $s$, 로그이면 $\log s$ (깊은 꼬리는 로그를 따로 계산)
\Else \Comment{큰 쪽 $=1-s\ge1/2$}
  \State \Return $1-s$, 로그이면 $\operatorname{log1p}(-s)$
\EndIf
\end{algorithmic}
\caption{작은 쪽 규칙 (stardust 분포 객체의 \code{cdf}, \code{ccdf}, \code{logcdf}, \code{logccdf})}
\label{alg:small-side}
\end{algorithm}

\cref{alg:small-side}에서 큰 쪽을 $1-s$로 만드는 뺄셈은 해롭지 않다. 결과가 $1/2$ 이상이므로 $s$의 상대오차가 결과에서는 $s/(1-s)\le1$배로 오히려 작아지기 때문이다. 큰 쪽의 로그는 $\operatorname{log1p}(-s)$로 계산한다. $\log(1-s)$로 계산하면 $s$가 작을 때 $1-s$에서 반올림이 생긴다. $s$가 $u/2$보다 작으면 큰 쪽은 반올림해서 그냥 $1$이다.

꼬리가 더 깊으면 확률 자체가 binary64로 표현되지 않는다. $1-\Phi(40)\approx10^{-350}$은 가장 작은 비정규 수 $4.9\times10^{-324}$보다 작아 $0$이 된다. 그러나 우도를 곱하거나 p값을 비교하는 계산에는 그런 확률도 필요하므로 확률의 로그를 확률을 거치지 않는 식으로 계산한다. 정규 분포의 \code{log\_ndtr}는 두 영역으로 나뉜다. $x\ge0$에서는 $\operatorname{log1p}(-\Phi(-x))$이고 $x<0$에서는 $\Phi(x)=\mathrm{erfc}(-x/\sqrt2)/2=e^{-x^2/2}\,\mathrm{erfcx}(-x/\sqrt2)/2$를 써서 $-(x/2)\cdot x+\log(\mathrm{erfcx}(-x/\sqrt2)/2)$로 계산한다($x/2$를 먼저 곱해 $|x|>1.3\times10^{154}$에서 넘치지 않게 한다). 척도를 곱한 여오차 함수 $\mathrm{erfcx}(t)=e^{t^2}\mathrm{erfc}(t)$는 Mills 비와 같은 양이고 확률을 거치지 않는다(85장). 감마와 베타는 불완전 감마·베타 함수의 앞 인자의 로그에서 얻는다(87--88장). 확률을 먼저 계산하고 로그를 취하면 $0$의 로그가 $-\infty$가 된다.

\subsection{Discrete Distribution Functions as Sums}

닫힌 꼴의 분포 함수가 없는 이산 분포(초기하, 베타이항 등)에서는 질량 함수를 더한다. 이때도 작은 쪽 규칙을 쓰며 정확도는 더하는 방식에 달려 있다. stardust는 꼬리의 첫 항을 $1$로 놓고 이웃한 항의 비를 정확한 정수 인자 넷의 몫 $(a_1a_2)/(b_1b_2)$로 만들어 double-double로 곱해 나간다(\cref{alg:tail-sum}). 질량 함수가 꼬리 쪽으로 줄어드는 방향으로만 가므로 6장에서 본 앞방향 점화식의 안정한 방향이다. 그래서 수천 항을 가도 로그 질량 함수로 다시 맞출 필요가 없다. 항이 합의 $2^{-60}$배 아래로 내려가면 남은 항을 마지막 항 $t$와 마지막 비 $q$의 등비급수 $tq/(1-q)$로 어림해 더한다. 비가 계속 줄어드므로 이 어림은 참값보다 조금 크다. 매끄럽고 아주 긴 꼬리는 Euler--Maclaurin 공식(3장)과 Gauss 구적으로 바꾸어 계산한다.

\begin{algorithm}[htbp]
\begin{algorithmic}[1]
\Require 시작점 $k$, 방향 $\pm1$, 끝, 비 $p(j\pm1)/p(j)$의 네 정수 인자
\Ensure $\sum_jp(j)/p(k)$
\State $t\gets1$ (double-double); $s\gets1$; $c\gets0$
\For{$j=k,k\pm1,\dots$ (많아야 $2^{20}$항)}
  \State $q\gets(a_1a_2)/(b_1b_2)$ (double-double); $t\gets t\cdot q$
  \State Neumaier 합으로 $s\gets s+t_{\mathrm{hi}}$, $c\gets c+t_{\mathrm{lo}}$
  \State \textbf{if} $t_{\mathrm{hi}}\le2^{-60}s$ \textbf{then} \Return $s+c+t_{\mathrm{hi}}q\min(1/(1-q),\text{남은 점 수})$
\EndFor
\State 끝에 닿지 못했으면 NaN \Comment{잘린 합은 내지 않는다}
\end{algorithmic}
\caption{비 점화식으로 더하는 이산 꼬리 (stardust \code{csrc/stats/discrete.hpp}의 \code{tail\_sum\_dd})}
\label{alg:tail-sum}
\end{algorithm}

\subsection{Measured Tails}

\cref{tab:tails}에서는 따로 빌드한 stardust를 십진 참값과 견주었다. \code{ccdf}는 $x=37$($\bar\Phi\approx6\times10^{-300}$)까지 상대오차가 $4.4\times10^{-16}$ 이하였고 같은 점에서 $1-$\code{cdf}는 $x=5$부터 자릿수를 잃어 $x=20$에서는 값이 $0$이었다. \code{logccdf}는 확률이 $0$으로 언더플로하는 $x=40$, $100$에서도 상대오차가 $2\times10^{-17}$, $4\times10^{-17}$이었다. \cref{fig:tails-curve}에 같은 비교를 $x$에 따라 그렸다.

\begin{figure}[htbp]
\centering
\begin{tikzpicture}
\begin{semilogyaxis}[sdplot,xlabel={$x$},ylabel={$\bar\Phi(x)$의 상대오차},xmin=0,xmax=38,ymin=1e-18,ymax=3,width=0.68\textwidth,height=0.42\textwidth]
\addplot[black,mark=*,mark size=0.7pt] table[x=x,y=ccdf]{figures/09/tails_curve.dat};
\addlegendentry{\code{ccdf}}
\addplot[black!50,thick] table[x=x,y=naive]{figures/09/tails_curve.dat};
\addlegendentry{$1-$\code{cdf}}
\end{semilogyaxis}
\end{tikzpicture}
\caption{표준 정규 위 꼬리의 상대오차(따로 빌드한 stardust; 참값은 Mills 비의 연분수를 십진 산술로, \code{figures/09/tails\_curve.py}). $1-$\code{cdf}는 $x\approx8.3$에서 값이 $0$이 되어 상대오차가 $1$에 붙는다. $10^{-18}$은 오차 $0$을 나타낸다. $x=38$에서는 $\bar\Phi\approx3\times10^{-317}$이 비정규 수라 상대오차가 $3\times10^{-9}$이고 계약은 이 영역에서 절대 $45\cdot2^{-1074}$이다.}
\label{fig:tails-curve}
\end{figure}

\begin{table}[htbp]
\begin{sdtable}{colspec={cccc}, row{1}={font=\small\bfseries}}
\toprule
$x$ & $\bar\Phi(x)$ & \code{ccdf}의 상대오차 & $1-$\code{cdf}의 상대오차 \\
\midrule
$3$ & $1.35\times10^{-3}$ & $4.4\times10^{-16}$ & $6.6\times10^{-15}$ \\
$5$ & $2.87\times10^{-7}$ & $1.9\times10^{-16}$ & $1.5\times10^{-10}$ \\
$8$ & $6.22\times10^{-16}$ & $1.2\times10^{-16}$ & $7.1\times10^{-2}$ \\
$20$ & $2.75\times10^{-89}$ & $3.6\times10^{-16}$ & $1$ (값이 $0$) \\
$37$ & $5.7\times10^{-300}$ & $3.0\times10^{-16}$ & $1$ (값이 $0$) \\
\bottomrule
\end{sdtable}
\caption{표준 정규 분포의 위 꼬리를 따로 빌드한 stardust의 \code{ccdf}와 $1-$\code{cdf}로 계산한 상대오차(참값은 Mills 비의 연분수를 십진 산술로; \code{figures/09/tails\_sd.py}).}
\label{tab:tails}
\end{table}

분위수 함수에도 같은 문제가 있어 위쪽 꼬리는 \code{iccdf}로, 로그 확률을 받는 역함수는 \code{ilogcdf}, \code{ilogccdf}로 따로 둔다(\cref{sec:computing-quantiles}). 이 네 짝과 그 계약이 106장의 분포 틀이 정한 공통 인터페이스이다.

\section{Testing Probability Code}
\label{sec:prob-testing}

확률 코드의 시험은 세 층으로 한다. 첫째는 정확한 참값이다. 이 장에서는 감마 분포 함수를 Poisson 합의 십진 산술로, 정규 꼬리를 Mills 비의 연분수로, Poisson 누적 확률을 닫힌 꼴로 계산했다. stardust 자신은 Arb의 구간 산술로 만든 참값 파일과 견주어 상대 $10^{-14}$의 계약(설계 문서 T10)을 확인하고 scipy와 비트까지 같다고는 약속하지 않는다(비트 정책). 둘째는 항등식이다. \code{cdf}$+$\code{ccdf}$=1$은 작은 쪽의 상대오차 안에서, \code{icdf(cdf(k))==k}는 이산 분포에서 정확히 성립해야 하고 \code{icdf}의 결과는 \cref{def:backward-quantile}의 계약을 만족하는지 이웃한 정수의 누적 확률로 확인할 수 있다. 셋째는 통계적 시험이다. 표본기는 정확한 분포를 내지 않고 그 분포를 따르는 열을 내므로 카이제곱이나 Kolmogorov--Smirnov 검정을 거쳐야 하고(7장 \cref{tab:fy-chi2}처럼) 그때 p값 자체도 \cref{thm:inverse-transform}에 따라 균등해야 한다. 같은 시험을 여러 번 돌려 p값의 분포가 균등한지 보는 것이 한 번의 통과보다 강한 증거이다.

\begin{table}[htbp]
\begin{sdtable}{colspec={X[1.6,l]X[2.6,l]X[2,l]}, row{1}={font=\small\bfseries}, cells={font=\small}}
\toprule
대상 & 무엇과 견주었나 & 결과 \\
\midrule
정규 \code{ccdf}, \code{logccdf} & Mills 비 연분수의 십진 값 & 상대 $4.4\times10^{-16}$ 이하($x\le37$), 로그 $4\times10^{-17}$ 이하 \\
Poisson \code{icdf}, \code{iccdf} & 작은 쪽 항등식, 정확한 누적 확률로 본 계약 & 모두 통과(\cref{tab:quantile-tests}) \\
감마 합의 정규 근사 & 십진 산술의 정확한 분포 함수 & $\sqrt n\cdot$오차 $\to0.1330$ \\
표본기 & 큰 수의 법칙, Glivenko--Cantelli의 속도 & $1/\sqrt n$ 띠 안(\cref{fig:lln,fig:gc}) \\
\code{Multivariate}\allowbreak\code{Normal.cdf} & 일차원 적분의 참값 & $300/300$이 보고 오차 안(\cref{fig:mvn-cover}) \\
\code{permutation\_test} & 귀무가설 아래 p값의 분포 & 초균등(\cref{fig:perm-p}) \\
\bottomrule
\end{sdtable}
\caption{이 장의 검증 요약(따로 빌드한 stardust; 스크립트는 \code{figures/09/}).}
\label{tab:prob-tests}
\end{table}

\section{APIs in This Chapter}
\label{sec:probability-apis}

\cref{tab:probability-apis}에 이 장의 이론을 구현한 stardust API를 모았다. 분포 틀 전체는 106장, 분포별 계산식은 107--108장, 난수와 표집은 104--105장에서 다룬다.

\begin{table}[htbp]
\begin{sdtable}{colspec={X[2.2,l]X[3,l]X[1.5,l]}, row{1}={font=\small\bfseries}}
\toprule
API & 이 장에서 다룬 것 & 설명한 곳 \\
\midrule
이산 분포의 \code{icdf}, \code{iccdf}, \code{ilogcdf}, \code{ilogccdf} & 후방 오차 계약, 갤로핑과 이분, 작은 쪽 비교, 원자에서의 정착 & \cref{alg:discrete-icdf,alg:settle}, \cref{tab:quantile-tests} \\
연속 분포의 \code{icdf}, \code{iccdf} & 분위수의 조건수, 로그 쪽 Newton & \cref{alg:normal-icdf} \\
\code{cdf}, \code{ccdf}, \code{logcdf}, \code{logccdf} & 작은 쪽 규칙, 깊은 로그 꼬리 & \cref{alg:small-side}, \cref{tab:tails} \\
분포 객체의 \code{sample} & 균등수 $(m+1/2)2^{-53}$, 역변환, 분포별 방법 & \cref{alg:by-inverse}, \cref{tab:samplers} \\
\code{random.Generator}, \code{spawn} & 카운터 기반 열과 독립 & \cref{alg:philox-words} \\
\code{htest.ks\_1samp} & 경험 분포 함수와 $D_n$, Kolmogorov 극한 & \cref{alg:ks-stat,alg:kolmogorov} \\
\code{htest.}\allowbreak\code{epps\_singleton\_2samp} & 경험 특성함수, 다변량 CLT & \cref{alg:epps} \\
\code{htest.permutation\_test}, \code{monte\_carlo\_test}, \code{bootstrap} & 큰 수의 법칙, p값 $(b+1)/(n+1)$ & \cref{prop:mc-pvalue}, \cref{alg:resample-p} \\
\code{MultivariateNormal.cdf}, \code{sample} & 정규 벡터의 조건부 분포, 순차 조건화, 준몬테카를로 오차 & \cref{alg:genz} \\
이산 분포의 \code{cdf}, \code{mean}, \code{skewness}, \code{kurtosis} & 비 점화식 꼬리 합, 중심 모멘트 합 & \cref{alg:tail-sum,alg:central-moments} \\
\bottomrule
\end{sdtable}
\caption{이 장의 이론을 쓰는 stardust API.}
\label{tab:probability-apis}
\end{table}

\section{Concluding Remarks}

전체 측도가 $1$인 측도론에 독립과 조건부 기댓값을 더하면 확률론이 된다. 분포는 누적분포함수로, 누적분포함수는 다시 분위수 함수로 표현된다. 분위수 함수가 균등 확률변수를 원하는 분포로 바꾼다는 정리에서 난수 생성이 나오고 검정은 그 역인 확률 적분 변환에 기댄다. 조건부 기댓값은 Radon--Nikodym 도함수로 정의되지만 $L^2$에서는 직교 사영이어서 회귀의 언어로 다시 읽힌다. 극한 정리는 표본평균이 참평균으로 가고 그 오차가 $\sigma/\sqrt n$의 정규 분포를 따른다고 말하며 실제 속도는 왜도가 정하는 Edgeworth의 첫 항으로 잘 설명된다. 구현에서는 수학적으로 같은 양이 binary64에서는 다른 값이었다. $1-F$와 $\bar F$, 계산한 확률의 로그와 직접 계산한 로그 확률이 그렇다. 이산 분위수는 조건수가 무한대라서 전방 오차가 아니라 후방 오차로만 약속할 수 있었다. 측도론 위의 확률론을 다룬 표준적인 참고서로는 Billingsley\cite{billingsley1995}와 Durrett\cite{durrett2019}가 있고 극한 정리의 정밀한 형태는 Petrov\cite{petrov1975}, 고전적인 직관은 Feller\cite{feller1968}에서 얻을 수 있다. 1권은 여기서 끝나고 2권에서는 컴퓨터 산술과 반올림 오차를 다룬다.

\chapterreferences
